\documentclass[12pt]{article}

% --- Packages ---
\usepackage[margin=1in]{geometry}
\usepackage{amsmath,amssymb,amsthm,mathtools}
\usepackage{enumitem}
\usepackage{graphicx}
\usepackage{hyperref}
\usepackage{natbib}
\usepackage{booktabs}
\usepackage{caption}
\usepackage{float}
\usepackage{xcolor}

% --- Theorem environments ---
\newtheorem{theorem}{Theorem}[section]
\newtheorem{proposition}[theorem]{Proposition}
\newtheorem{lemma}[theorem]{Lemma}
\newtheorem{corollary}[theorem]{Corollary}
\newtheorem{conjecture}[theorem]{Conjecture}
\theoremstyle{remark}
\newtheorem{remark}[theorem]{Remark}
\theoremstyle{definition}
\newtheorem{definition}[theorem]{Definition}
\newtheorem{assumption}[theorem]{Assumption}
\newtheorem{example}[theorem]{Example}

% --- Shortcuts ---
\newcommand{\E}{\mathbb{E}}
\newcommand{\Var}{\mathrm{Var}}
\newcommand{\Cov}{\mathrm{Cov}}
\newcommand{\Prob}{\mathbb{P}}
\newcommand{\R}{\mathbb{R}}
\newcommand{\N}{\mathbb{N}}
\newcommand{\FI}{\mathrm{FI}}
\newcommand{\PI}{\mathrm{PI}}
\newcommand{\What}{\widehat{W}}
\newcommand{\Ftrue}{\Prob_{\mu_1}}
\newcommand{\Fbad}{\Prob_{\mu_2}}
\DeclareMathOperator{\sgn}{sgn}

\title{First Passage Times in Mean-Variance Games:\\
The Risk Reversal of Partial Information}

\author{%
  Author Name\thanks{School of Mathematical Sciences, Fudan University.
  Email: \texttt{email@fudan.edu.cn}.}
}

\date{April 2026}

\begin{document}

\maketitle

\begin{abstract}
We study first passage times in symmetric mean-variance portfolio games
under full and partial information, extending the framework of
\citet{HuangSun2025}.
Our main finding is a \emph{risk reversal}: the direction of the
first passage comparison between full information~(FI) and partial
information~(PI) investors depends on the unobserved state.
Under the good state ($\mu = \mu_1$), the PI investor is dramatically
safer---her ruin probability is bounded above by the FI infinite-horizon
ruin probability. Our theoretical lower bound gives a Lundberg exponent
ratio of $\mu_1/\theta(p) \approx 1.82$; Monte Carlo estimates suggest
the true ratio is $5$--$6$ times, with the gap attributable to the
hedging component of the PI strategy.
Under the bad state ($\mu = \mu_2$), the comparison reverses: the PI
investor over-invests by a factor of up to $5.5\times$, amplifying ruin
probability by up to $300\times$ in the infinite-horizon limit.
The sensitivity of ruin probability to competition intensity~$\lambda$
is $5$--$7$ times greater under partial information, providing a precise
characterization of the ``downward amplification'' conjectured by
\citet{HuangSun2025}.
We establish these results through four theorems---an exponential
supermartingale bound, a Dubins--Schwarz quadratic variation comparison
under the martingale measure, a bad-state amplification inequality via
$L^2$ coupling, and a $\lambda$-sensitivity comparison---complemented
by comprehensive Monte Carlo analysis and a perturbation expansion for
heterogeneous investors.
\end{abstract}

\bigskip
\noindent\textbf{Keywords:}
mean-variance games, partial information, first passage time,
ruin probability, Nash equilibrium, relative performance, Bayesian filtering,
systemic risk

\bigskip
\noindent\textbf{JEL Classification:} C73, D82, G11

\bigskip
\noindent\textbf{MSC 2020:} 91A15, 91G10, 60G35, 60G40

\newpage

%% ========================================================================
\section{Introduction}\label{sec:intro}
%% ========================================================================

Competition among fund managers is a pervasive feature of modern
financial markets.
Institutional investors are evaluated not only on absolute returns but
also on performance relative to their peers, creating strong incentives
for benchmark tracking and herd behavior.
At the same time, many of the fundamental parameters governing
asset returns---expected growth rates, regime probabilities, structural
break points---are not directly observable.
Fund managers must learn about these parameters in real time from market
data, making investment decisions under partial information.
The intersection of these two forces---competitive pressure and
informational uncertainty---gives rise to subtle effects on tail risk
that are the subject of this paper.

The mean-variance game framework provides a natural setting for
studying competitive portfolio choice.
\citet{BasakMakarov2014} pioneered the study of two-fund
continuous-time competition with relative performance concerns,
showing that benchmarking induces risk-shifting and excess volatility.
\citet{EspinosaTouzi2015} analyzed mean-field games with relative utility
in the large-population limit,
and \citet{LackerZariphopoulou2019} established existence results
for $N$-player forward performance games under general preferences.
In the related literature on executive stock options under partial
information, \citet{Henderson2020} studied the interplay between
learning and exercise decisions, demonstrating that informational
frictions can alter the qualitative structure of optimal policies.
The stochastic ordering approach of \citet{BaeuerleBayraktar2013}
provides a general framework for comparing stochastic control problems
under different information structures.
In the CRRA setting, \citet{Deng2022} studied relative wealth games
under partial information with heterogeneous priors, showing that
competitive pressure induces a herd effect whereby investors trade
more aggressively; \citet{ZhangYan2024} extended this to a factor
model and mean-field limit.
These works establish that competition amplifies the \emph{magnitude}
of trading activity, but do not study how the interaction of
competition and partial information affects \emph{first passage times}
or ruin probabilities---the direction we pursue here.
On the systemic risk side, \citet{Carmona2018} modeled inter-bank
lending with delay, showing that strategic interaction can create
correlated default cascades; see also \citet{Leobacher2016} for
Bayesian ruin theory under partial information in the single-agent
dividend optimization context.

\citet{HuangSun2025} brought competitive games and partial information
together by solving $N$-player mean-variance games under both full and
partial information in a Bayesian filtering framework.
Their model yields explicit Nash equilibria in all four settings
(constant or alternating drift, full or partial information),
and their numerical experiments reveal a striking phenomenon:
under partial information, relative performance concerns
\emph{amplify downward risk} in a way that is absent under full
information.
Specifically, the probability that a large fraction of investors
simultaneously fall below a warning threshold increases dramatically
when investors compete under partial information, but remains
essentially unchanged under full information.
\citet{HuangSun2025} identify this as an open problem:
\begin{quote}
``It is of interest as future research to investigate if the
`downward amplification' effect of partial information on investors'
wealth can be theoretically proved and characterized, beyond the
motivating numerical illustration in Figure~4.''
\end{quote}

The present paper addresses this open problem and arrives at a
conclusion that is both more precise and more surprising than the
original conjecture suggests.
We show that the first passage time comparison between FI and PI
investors exhibits a \emph{risk reversal} that depends on the
unobserved true state of the drift:

\begin{enumerate}[label=(\roman*)]
\item \textbf{Good-state safety.}
  When the true drift is high ($\mu = \mu_1$), the PI investor is
  dramatically \emph{safer} than the FI investor.
  The PI investor, not knowing the state is favorable, under-invests
  relative to her FI counterpart.
  This conservative behavior reduces portfolio volatility more than
  it reduces expected returns (the Sharpe ratio is preserved), leading
  to an exponentially smaller ruin probability.
  We prove that the PI finite-horizon ruin probability is bounded by
  the FI \emph{infinite-horizon} ruin probability
  (Theorem~\ref{thm:expo_bound}), and we establish a comparison under
  the martingale measure via quadratic variation ordering
  (Theorem~\ref{thm:Q_comparison}).
  Monte Carlo simulations confirm that the effective Lundberg exponent
  of the PI investor is $5$--$6$ times that of the FI investor across
  all competition levels.

\item \textbf{Bad-state amplification.}
  When the true drift is low ($\mu = \mu_2$), the comparison reverses.
  The PI investor, unaware that the state is unfavorable, over-invests
  by a factor of up to $\theta(p)/\mu_2$ (which equals $5.5$ at the
  \citeauthor{HuangSun2025} parameters).
  This over-investment amplifies ruin probability by up to $300\times$
  in the infinite-horizon limit (Theorem~\ref{thm:bad_state}).
  The bad-state amplification is \emph{exponentially larger} than the
  good-state dampening, explaining why the mixture (unconditional)
  comparison is dominated by the bad-state effect at moderate priors.

\item \textbf{Sensitivity amplification.}
  The log-ruin probability is $5$--$7$ times more sensitive to the
  competition parameter~$\lambda$ under PI than under FI
  (Theorem~\ref{thm:sensitivity}).
  This is the ``downward amplification'' of \citet{HuangSun2025}:
  increasing $\lambda$ erodes the PI safety margin proportionally
  faster, even though the PI ruin level is absolutely lower under
  the good state.
\end{enumerate}

Together, results (i)--(iii) reveal a ``paradox of ignorance'':
not knowing the true state \emph{protects} the investor when the state
is good (because conservatism reduces volatility) and \emph{endangers}
her when the state is bad (because misplaced optimism amplifies
exposure).
The net effect depends on the prior and on the magnitude of the drift
gap $\delta := \mu_1 - \mu_2$.

\medskip
\noindent\textbf{Contributions.}
Our specific contributions are:
\begin{enumerate}[label=\arabic*.]
\item Four rigorous theorems characterizing the first passage
  comparison: an exponential supermartingale bound
  (Theorem~\ref{thm:expo_bound}), a Dubins--Schwarz comparison under
  the martingale measure (Theorem~\ref{thm:Q_comparison}),
  a bad-state ruin amplification bound via $L^2$ coupling
  (Theorem~\ref{thm:bad_state}), and a $\lambda$-sensitivity
  comparison (Theorem~\ref{thm:sensitivity}).
\item A comprehensive numerical analysis across the full
  $(D, \lambda)$ parameter space, confirming that the PI investor is
  uniformly safer under the good state and identifying the bad-state
  reversal region.
\item A perturbation expansion for heterogeneous investors
  ($\gamma_A \neq \gamma_B$), showing that heterogeneity breaks the
  $\lambda$-independence of the PI strategy and reduces co-drawdown
  through investor decorrelation
  (Proposition~\ref{prop:heterogeneous}).
\item An honest treatment of the Girsanov obstruction
  (Remark~\ref{rem:girsanov}), explaining precisely why the
  martingale-measure comparison does not transfer directly to the
  physical measure and formulating this as an open problem.
\end{enumerate}

\medskip
\noindent\textbf{Organization.}
Section~\ref{sec:model} presents the model and recalls the Nash
equilibria.
Section~\ref{sec:results} states the main results.
Section~\ref{sec:proofs} contains proofs.
Section~5 provides numerical analysis, and
Section~6 discusses extensions and open problems.

%% ========================================================================
\section{Model}\label{sec:model}
%% ========================================================================

We adopt the framework of \citet{HuangSun2025} and work throughout
under the \emph{constant-drift, symmetric-investor} specification.

\subsection{Financial market}\label{sec:market}

Let $(\Omega, \mathcal{F}, \{\mathcal{F}_t\}_{t \ge 0}, \Prob)$ be a
filtered probability space satisfying the usual conditions, supporting
a standard Brownian motion~$W$.
A single risky asset (stock) has price process
\begin{equation}\label{eq:stock}
  \frac{dS(t)}{S(t)} = \mu\, dt + \sigma\, dW(t), \qquad S(0) = s_0 > 0,
\end{equation}
where $\sigma > 0$ is known and $\mu \in \{\mu_1, \mu_2\}$ is a fixed but
possibly unknown constant with $\mu_1 > \mu_2 > 0$.
The risk-free rate is $r = 0$.
We write $\delta := \mu_1 - \mu_2 > 0$ for the drift gap.

Two probability measures play a central role.
When the true drift is $\mu_k$, we write~$\Prob_{\mu_k}$ for the
corresponding probability measure.
Under $\Prob_{\mu_k}$, the process $W$ is a standard Brownian motion.

\subsection{Investors and objectives}\label{sec:investors}

There are $N \ge 2$ investors, each choosing a dollar investment
$\pi_i(t)$ in the stock at time~$t$.
Investor~$i$'s wealth satisfies
\begin{equation}\label{eq:wealth}
  dX_i(t) = \pi_i(t)\,\mu\,dt + \pi_i(t)\,\sigma\,dW(t),
  \qquad X_i(0) = x > 0.
\end{equation}

\begin{assumption}[Symmetry]\label{ass:sym}
All investors share the same risk aversion $\gamma > 0$,
the same relative performance weight
$\lambda^M = \lambda^V = \lambda \in [0,1)$,
and the same initial wealth~$x$ and warning threshold~$D < x$.
\end{assumption}

Under Assumption~\ref{ass:sym}, the constants defined in
\citet[Eq.~(3.6)--(3.7)]{HuangSun2025} reduce to
\begin{equation}\label{eq:kappa}
  \kappa_i = \kappa := \frac{1}{\gamma}\,,
  \qquad
  \bar{\kappa} = \frac{1}{\gamma}\,,
  \qquad
  \bar{\lambda}^V = \lambda.
\end{equation}
The effective risk aversion for the game becomes $\gamma(1-\lambda)$
(see \citet[Remark~3.4]{HuangSun2025}).

Each investor maximizes the mean-variance objective with relative
performance:
\begin{equation}\label{eq:objective}
  J_i = \E\bigl[X_i(T) - \lambda\,\bar{X}(T)\bigr]
      - \frac{\gamma}{2}\,\Var\bigl[X_i(T) - \lambda\,\bar{X}(T)\bigr],
\end{equation}
where $\bar{X} := \frac{1}{N}\sum_{j=1}^N X_j$ is the average wealth.

\subsection{Information regimes}\label{sec:info}

\paragraph{Full information (FI).}
All investors observe $\mu$ directly.
The relevant filtration is $\mathcal{F}^{\mu,S}$, and investors know
the true value of the drift.

\paragraph{Partial information (PI).}
Investors observe only the stock price $S$; the drift
$\mu \in \{\mu_1, \mu_2\}$ is unknown.
The prior probability is
$p := \Prob(\mu = \mu_1 \mid \mathcal{F}_0) \in (0,1)$.
Investors use the filtration $\mathcal{F}^S$ generated by~$S$.

\subsection{Nash equilibria}\label{sec:equil}

We recall the equilibrium strategies from \citet{HuangSun2025},
specialized to the symmetric case with $r = 0$.

\subsubsection{Full information equilibrium}

By \citet[Theorem~3.1]{HuangSun2025}, when the true drift is
$\mu = \mu_k$, the FI Nash equilibrium strategy is
\begin{equation}\label{eq:piFI}
  \pi_k^{\FI} := \frac{\mu_k}{\sigma^2\,\gamma(1-\lambda)}\,,
\end{equation}
a deterministic constant.
The corresponding wealth process under $\Prob_{\mu_k}$ is
\begin{equation}\label{eq:XFI}
  X_k^{\FI}(t) = x + \pi_k^{\FI}\,\mu_k\,t
              + \pi_k^{\FI}\,\sigma\,W(t),
  \qquad t \in [0,T].
\end{equation}
This is Brownian motion with constant drift
$\pi_k^{\FI}\mu_k = \mu_k^2/[\sigma^2\gamma(1-\lambda)]$ and constant
diffusion coefficient $\pi_k^{\FI}\sigma = \mu_k/[\sigma\gamma(1-\lambda)]$.

\begin{remark}\label{rem:FI_state}
The FI investor \emph{knows} the true state and adapts her strategy
accordingly.
Under the good state, she invests
$\pi_1^{\FI} = \mu_1/[\sigma^2\gamma(1-\lambda)]$; under the bad state,
$\pi_2^{\FI} = \mu_2/[\sigma^2\gamma(1-\lambda)]$.
The ratio $\pi_1^{\FI}/\pi_2^{\FI} = \mu_1/\mu_2$ can be large
(it equals $10$ at the \citeauthor{HuangSun2025} parameters
$\mu_1 = 0.2$, $\mu_2 = 0.02$).
\end{remark}

\subsubsection{Partial information equilibrium}

Under PI, the posterior probability
$P(t) := \Prob(\mu = \mu_1 \mid \mathcal{F}^S_t)$
satisfies the filtering SDE
(\citealt[Lemma~3.2]{HuangSun2025}):
\begin{equation}\label{eq:filter}
  dP(t) = \beta\bigl(P(t)\bigr)\,d\What(t),
  \qquad P(0) = p,
\end{equation}
where the filter volatility and estimated drift are
\begin{equation}\label{eq:theta_beta}
  \theta(q) := \delta\,q + \mu_2,
  \qquad
  \beta(q) := \frac{\delta}{\sigma}\,q(1-q),
\end{equation}
and $\What$ is the innovation Brownian motion, a standard Brownian motion
under $\mathcal{F}^S$.

By \citet[Theorem~3.2, Remark~3.7]{HuangSun2025}, the PI Nash
equilibrium strategy in the symmetric case with $r = 0$ is
\begin{equation}\label{eq:piPI}
  \pi^*(t,q) = \frac{g(t,q)}{\gamma(1-\lambda)}\,,
\end{equation}
where $g\colon [0,T] \times [0,1] \to \R$ encodes both the myopic and
hedging components:
\begin{equation}\label{eq:g_decomp}
  g(t,q)
  = \underbrace{\frac{\theta(q)}{\sigma^2}}_{\text{myopic}}
  \;-\; \underbrace{\frac{\beta(q)}{\sigma}\,
    \frac{\partial_q c(t,q)}{\kappa}}_{\text{hedging}}\,,
\end{equation}
with $c(t,q)$ solving the Cauchy problem
(\citealt[Lemma~3.3, Eq.~(3.27)]{HuangSun2025})
\begin{equation}\label{eq:cauchy}
  \partial_t c + \frac{\beta(q)^2}{2}\,\partial_{qq} c
  - \frac{\beta(q)\,\theta(q)}{\sigma}\,\partial_q c
  + \kappa\,\frac{\theta(q)^2}{\sigma^2} = 0,
  \qquad c(T,q) = 0.
\end{equation}

\begin{remark}[Structural property: $g$ is $\lambda$-independent]
\label{rem:g_indep}
A key feature of the symmetric equilibrium is that $\kappa = 1/\gamma$,
$\bar{\lambda}^V = \lambda$, and the competition-induced modifications to
the Cauchy problem~\eqref{eq:cauchy} cancel exactly
(\citealt[Remark~3.7]{HuangSun2025}).
As a result, $g(t,q)$ does not depend on~$\lambda$, and the equilibrium
strategy $\pi^*(t,q) = g(t,q)/[\gamma(1-\lambda)]$ depends on~$\lambda$
only through the explicit prefactor $1/(1-\lambda)$.
This property is central to our sensitivity analysis
(Theorem~\ref{thm:sensitivity}) and breaks down under heterogeneous
risk aversion (Proposition~\ref{prop:heterogeneous}).
\end{remark}

\subsection{Wealth processes under both measures}
\label{sec:wealth_both}

A key aspect of our analysis is that the \emph{same} PI strategy is
deployed regardless of the true state, while the FI strategy adapts to
the state.
We now make this precise.

Under the true measure $\Prob_{\mu_k}$ (where $W$ is standard Brownian
motion), the PI wealth process satisfies
\begin{equation}\label{eq:XPI}
  dX^{\PI}(t) = \pi^*\bigl(t, P(t)\bigr)\,\mu_k\,dt
              + \pi^*\bigl(t, P(t)\bigr)\,\sigma\,dW(t).
\end{equation}
Importantly, the filter $P(t)$ evolves according to
different dynamics under~$\Prob_{\mu_1}$ and~$\Prob_{\mu_2}$.
Under $\Prob_{\mu_1}$, the innovation process is
$d\What(t) = dW(t) + [\mu_1 - \theta(P(t))]/\sigma\,dt$,
and $P(t)$ is a submartingale converging to~$1$.
Under $\Prob_{\mu_2}$, the innovation is
$d\What(t) = dW(t) + [\mu_2 - \theta(P(t))]/\sigma\,dt$,
and $P(t)$ is a supermartingale converging to~$0$.

\begin{remark}[Investment levels under each state]
\label{rem:investment}
Under $\Prob_{\mu_1}$, the PI investor's expected drift estimate
satisfies $\E_{\mu_1}[\theta(P(t))] < \mu_1$ for all $t > 0$ when
$p < 1$.
The PI investor therefore \emph{under-invests} relative to the FI
strategy $\pi_1^{\FI}$.

Under $\Prob_{\mu_2}$, the PI investor's expected drift estimate
satisfies $\E_{\mu_2}[\theta(P(t))] > \mu_2$ for all $t > 0$ when
$p > 0$.
The PI investor therefore \emph{over-invests} relative to the FI
strategy $\pi_2^{\FI}$.
At time zero, the over-investment ratio is
\begin{equation}\label{eq:over_invest}
  \frac{\pi^*(0)}{\pi_2^{\FI}} = \frac{\theta(p)}{\mu_2}
  = 1 + \frac{\delta p}{\mu_2}\,,
\end{equation}
which equals $5.5$ at the \citeauthor{HuangSun2025} parameters
($\delta = 0.18$, $p = 0.5$, $\mu_2 = 0.02$).
\end{remark}

\subsection{Definitions}\label{sec:definitions}

\begin{definition}[First passage time]\label{def:fpt}
For a wealth process $X$ with initial value $x > D$, the
\emph{first passage time} (ruin time) to the warning threshold~$D$ is
\begin{equation}\label{eq:fpt}
  \tau := \inf\{t \ge 0 : X(t) \le D\}.
\end{equation}
We write $\tau_k^{\FI}$ for the FI ruin time under state~$k$, and
$\tau^{\PI}$ for the PI ruin time.
The \emph{safety margin} is $a := x - D > 0$.
\end{definition}

\begin{definition}[Co-drawdown probability]\label{def:codrawdown}
Under Assumption~\ref{ass:sym}, all investors in the same information
regime are identical.
The co-drawdown probability for~$k$ of~$N$ investors under information
regime~$I \in \{\FI, \PI\}$ reduces to
$\Prob(\tau^I \le T)$.
\end{definition}

\begin{definition}[Expected wealth gap]\label{def:gap}
The expected wealth gap at horizon~$T$ under the true measure
$\Prob_{\mu_1}$ is
\begin{equation}\label{eq:gap_def}
  G(\lambda, T) := \E_{\mu_1}\bigl[X^{\FI}(T)\bigr]
                 - \E_{\mu_1}\bigl[X^{\PI}(T)\bigr].
\end{equation}
\end{definition}

%% ========================================================================
\section{Main Results}\label{sec:results}
%% ========================================================================

Throughout this section, Assumption~\ref{ass:sym} is in force with
$r = 0$, $\mu_1 > \mu_2 > 0$, and non-degenerate prior $p \in (0,1)$.
We write $a = x - D > 0$ for the safety margin.
We work with the myopic PI strategy
$\pi^*(t) = \theta(P(t))/[\sigma^2\gamma(1-\lambda)]$;
the results extend to the full strategy
$\pi^*(t) = g(t,P(t))/[\gamma(1-\lambda)]$ with the hedging
component, since the hedging demand is non-negative
(see Remark~\ref{rem:hedging_ext}).

\subsection{Expected wealth comparison (recap)}
\label{sec:wealth_comp}

We begin by recalling the expected wealth comparison from
\citet{HuangSun2025}, which motivates the first passage analysis.

\begin{theorem}[Suboptimality of partial information]
\label{thm:wealth}
Under $\Prob_{\mu_1}$ with $p \in (0,1)$:
\begin{enumerate}[label=(\roman*)]
\item The expected terminal wealth under PI is strictly less than
  under FI:
  $\E_{\mu_1}[X^{\PI}(T)] < \E_{\mu_1}[X^{\FI}(T)]$.
\item The expected wealth gap has the integral representation
  \begin{equation}\label{eq:gap_repr}
    G(\lambda, T) = \frac{\mu_1}{\gamma(1-\lambda)}
      \int_0^T \biggl[\frac{\mu_1}{\sigma^2}
      - \E_{\mu_1}\bigl[g\bigl(t, P(t)\bigr)\bigr]\biggr]\,dt.
  \end{equation}
\item The gap is strictly increasing in $\lambda$:
  $\partial G / \partial\lambda = G(\lambda,T)/(1-\lambda) > 0$.
\end{enumerate}
\end{theorem}

\begin{remark}\label{rem:gap_message}
Theorem~\ref{thm:wealth} shows that PI investors earn less \emph{on
average}.
However, the expected wealth comparison does not address the tail
behavior.
The first passage times in the remainder of this section show that
the \emph{direction} of the tail comparison reverses between the
two states, a phenomenon invisible to the expected wealth analysis.
\end{remark}

\subsection{Good-state first passage safety}
\label{sec:good_state}

The next three results establish that the PI investor has a smaller
ruin probability under the good state $\Prob_{\mu_1}$.

\begin{theorem}[Exponential safety bound]
\label{thm:expo_bound}
Under $\Prob_{\mu_1}$, for any $T > 0$:
\begin{equation}\label{eq:expo_bound}
  \Prob_{\mu_1}\bigl(\tau^{\PI} \le T\bigr)
  \;\le\;
  \exp\!\biggl(-\frac{2\mu_1\,a}{\pi_1^{\FI}\,\sigma^2}\biggr)
  \;=\;
  \Prob_{\mu_1}\bigl(\tau_1^{\FI} < \infty\bigr).
\end{equation}
That is, the PI finite-horizon ruin probability is bounded above by the
FI \emph{infinite-horizon} ruin probability.
\end{theorem}

\begin{remark}\label{rem:expo_interpret}
Since $\Prob_{\mu_1}(\tau_1^{\FI} < \infty)
  \ge \Prob_{\mu_1}(\tau_1^{\FI} \le T)$ trivially,
Theorem~\ref{thm:expo_bound} does not by itself establish
$\Prob_{\mu_1}(\tau^{\PI} \le T) \le \Prob_{\mu_1}(\tau_1^{\FI} \le T)$.
The gap between the finite-horizon and infinite-horizon FI ruin
probabilities is
$\Prob_{\mu_1}(\tau_1^{\FI} < \infty) - \Prob_{\mu_1}(\tau_1^{\FI} \le T)
  = e^{-2\gamma(1-\lambda)a}\,
  2\Phi\bigl(-a\sqrt{2\gamma(1-\lambda)/T}\,\bigr)$,
which is exponentially small in $a$ for fixed~$T$ and vanishes as
$T \to \infty$.
The full finite-horizon comparison
$\Prob_{\mu_1}(\tau^{\PI} \le T) \le \Prob_{\mu_1}(\tau_1^{\FI} \le T)$
is confirmed by extensive Monte Carlo simulation (Section~5) and
established under the martingale measure in
Theorem~\ref{thm:Q_comparison} below, but the transfer from the
martingale measure to the physical measure encounters a technical
obstruction (Remark~\ref{rem:girsanov}).
\end{remark}

\begin{remark}[Lundberg exponent interpretation]
\label{rem:lundberg}
The FI ruin probability satisfies
$\Prob_{\mu_1}(\tau_1^{\FI} < \infty) = e^{-\psi^{\FI} a}$
with Lundberg exponent
$\psi^{\FI} = 2\mu_1/(\pi_1^{\FI}\sigma^2) = 2\gamma(1-\lambda)$.
Theorem~\ref{thm:expo_bound} implies that the PI ruin probability
decays at least as fast:
$\Prob_{\mu_1}(\tau^{\PI} \le T) \le e^{-\psi^{\FI} a}$.
Monte Carlo estimates of the effective PI exponent yield
$\psi^{\PI} \approx 5$--$6 \times \psi^{\FI}$
(Table~\ref{tab:lundberg} in Section~5), reflecting the
exponentially stronger safety under partial information.
\end{remark}

\begin{theorem}[Martingale measure comparison]
\label{thm:Q_comparison}
Define $Q$ by the Radon--Nikodym density
\begin{equation}\label{eq:Q_density}
  \frac{dQ}{d\Prob_{\mu_1}}\bigg|_{\mathcal{F}_T}
  = \exp\!\biggl(-\frac{\mu_1}{\sigma}\,W(T)
    - \frac{\mu_1^2}{2\sigma^2}\,T\biggr).
\end{equation}
Under $Q$, both FI and PI wealth processes are continuous local
martingales, and $\widetilde{W}(t) := W(t) + (\mu_1/\sigma)\,t$ is a
standard $Q$-Brownian motion.
Then:
\begin{equation}\label{eq:Q_comp}
  Q\bigl(\tau^{\PI} \le T\bigr) \;\le\; Q\bigl(\tau_1^{\FI} \le T\bigr),
\end{equation}
with equality if and only if
$\pi^*(t, P(t)) = \pi_1^{\FI}$ for Lebesgue-a.e.\ $t \in [0,T]$,
$Q$-a.s.
\end{theorem}

\begin{remark}\label{rem:Q_intuition}
Under $Q$, both wealth processes are driftless: the FI process is a
scaled Brownian motion with quadratic variation
$U(T) = (\pi_1^{\FI}\sigma)^2 T$, while the PI process has random
quadratic variation
$V(T) = \sigma^2\int_0^T [\pi^*(s)]^2\,ds \le U(T)$ a.s.
The smaller quadratic variation implies fewer barrier crossings by
a Dubins--Schwarz time-change argument (see Section~\ref{sec:proof_Q}).
\end{remark}

\begin{corollary}[Constant damping]
\label{cor:constant}
If $\pi^*(t) \equiv c$ for a constant $c \in (0, \pi_1^{\FI})$, then
\begin{equation}\label{eq:const_comp}
  \Prob_{\mu_k}\bigl(\tau^{\PI} \le T\bigr)
  \;\le\;
  \Prob_{\mu_k}\bigl(\tau_1^{\FI} \le T\bigr)
\end{equation}
for $k = 1$, for all $T > 0$, under \emph{any} probability measure
equivalent to $\Prob_{\mu_1}$.
\end{corollary}

\begin{remark}\label{rem:constant_intuition}
Corollary~\ref{cor:constant} is elementary but illuminating.
When $c < \pi_1^{\FI}$, the PI wealth process is a deterministic
rescaling: $X^{\PI}(t) - x = (c/\pi_1^{\FI})(X_1^{\FI}(t) - x)$.
The PI process hits the barrier~$D$ if and only if the FI process
hits the lower barrier
$x - a\pi_1^{\FI}/c < D$, which is a strictly rarer event.
The general (time-varying) case of Theorem~\ref{thm:Q_comparison}
requires the Dubins--Schwarz reduction because the path-dependent
rescaling does not admit such a clean coupling.
\end{remark}

\begin{remark}[The Girsanov obstruction]
\label{rem:girsanov}
Theorem~\ref{thm:Q_comparison} establishes
$Q(\tau^{\PI} \le T) \le Q(\tau_1^{\FI} \le T)$ but does not immediately
yield the corresponding inequality under $\Prob_{\mu_1}$.
The Girsanov back-transformation gives
$\Prob_{\mu_1}(\tau^I \le T)
  = \E_Q[L(T)\,\mathbf{1}_{\{\tau^I \le T\}}]$,
where $L(T) = \exp[(\mu_1/\sigma)\widetilde{W}(T)
  - \mu_1^2 T/(2\sigma^2)]$ is the Radon--Nikodym density
$d\Prob_{\mu_1}/dQ\big|_{\mathcal{F}_T}$.
The density $L(T)$ depends on $\widetilde{W}(T)$, which correlates
with the ruin events $\{\tau^I \le T\}$ through the stock path.
Since $L(T)$ may correlate differently with $\{\tau^{\PI} \le T\}$ and
$\{\tau_1^{\FI} \le T\}$, the $Q$-ordering does not transfer to
$\Prob_{\mu_1}$ in general.

To see the difficulty concretely: on the event $\{\tau^I \le T\}$, the
wealth process has visited $D < x$, which requires a negative excursion
of the stock.
But $L(T)$ is increasing in $\widetilde{W}(T)$, so large values of
$L(T)$ are associated with paths that have a large \emph{terminal}
value of the stock---which are precisely the paths \emph{less likely}
to have experienced ruin.
The two events create opposite selection effects, and their interaction
depends on the portfolio strategy in a non-trivial way.

We leave the full physical-measure comparison as an open problem:
\begin{conjecture}\label{conj:physical}
Under the hypotheses of Theorem~\ref{thm:Q_comparison},
$\Prob_{\mu_1}(\tau^{\PI} \le T) \le \Prob_{\mu_1}(\tau_1^{\FI} \le T)$
for all $T > 0$.
\end{conjecture}
This conjecture is supported by $200{,}000$-path Monte Carlo
simulations across the full $(D, \lambda)$ parameter space, where the
ratio $\Prob_{\mu_1}(\tau^{\PI} \le T)/\Prob_{\mu_1}(\tau_1^{\FI} \le T)$
is uniformly less than~$1$ (Section~\ref{sec:num_phase}).
\end{remark}

\begin{remark}[Extension to the full strategy with hedging]
\label{rem:hedging_ext}
Theorems~\ref{thm:expo_bound} and~\ref{thm:Q_comparison} are stated
for the myopic PI strategy
$\pi^*(t) = \theta(P(t))/[\sigma^2\gamma(1-\lambda)]$.
They extend to the full strategy
$\pi^*(t) = g(t,P(t))/[\gamma(1-\lambda)]$ provided the hedging demand
$h(t,q) := \theta(q)/\sigma^2 - g(t,q) \ge 0$, which ensures
$\pi^{\text{full}}(t) \le \pi^{\text{myopic}}(t) \le \pi_1^{\FI}$.
The sign $h(t,q) \ge 0$ follows from the monotonicity
$\partial_q c(t,q) \ge 0$ of the value function with respect to the
posterior probability, which can be verified via a comparison principle
for the Cauchy problem~\eqref{eq:cauchy}
\citep[cf.][Remark~3.7]{HuangSun2025}.
Both proofs use only the upper bound $\pi^*(t) \le \pi_1^{\FI}$, which
holds a fortiori for the full strategy.
\end{remark}

\subsection{Bad-state amplification}
\label{sec:bad_state}

We now turn to the bad state $\Prob_{\mu_2}$, where the comparison
reverses.

\begin{proposition}[Over-investment under the bad state]
\label{prop:over_invest}
Under $\Prob_{\mu_2}$, for any $p \in (0,1)$ and $t > 0$:
\begin{equation}\label{eq:over_invest_exp}
  \E_{\mu_2}\bigl[\pi^*(t)\bigr] > \pi_2^{\FI}.
\end{equation}
Moreover, the initial over-investment ratio is
$\pi^*(0)/\pi_2^{\FI} = 1 + \delta p/\mu_2$, and the expected PI
portfolio converges to $\pi_2^{\FI}$ as $t \to \infty$.
\end{proposition}

\begin{theorem}[Bad-state ruin comparison]
\label{thm:bad_state}
Under $\Prob_{\mu_2}$, let
$\bar{\pi} := \pi^*(0) = \theta(p)/[\sigma^2\gamma(1-\lambda)]$.
Then:

\begin{enumerate}[label=(\roman*)]
\item (Infinite-horizon comparison.)
  For all $a > 0$ and $p \in (0,1)$:
  \begin{equation}\label{eq:bad_inf}
    \Prob_{\mu_2}\bigl(\tau^{\PI} < \infty\bigr)
    > \Prob_{\mu_2}\bigl(\tau_2^{\FI} < \infty\bigr).
  \end{equation}

\item (Quantitative bound.)
  For any $\epsilon \in (0, \delta p/\mu_2)$ and $h > 0$:
  \begin{equation}\label{eq:bad_quant}
    \Prob_{\mu_2}\bigl(\tau^{\PI} \le h\bigr)
    \;\ge\;
    \mathfrak{p}\!\bigl((1+\epsilon)a,\; \bar{\pi}\mu_2,\;
      \bar{\pi}\sigma,\; h\bigr)
    - \frac{K\,h^2}{\epsilon^2\,a^2}\,,
  \end{equation}
  where $\mathfrak{p}(a', b, c, T)$ denotes the Bachelier--L\'evy ruin
  probability for a Brownian motion with drift~$b$ and diffusion~$c$
  to reach level~$-a'$ by time~$T$, and
  $K = K(\delta, \sigma, \gamma, \lambda, p)$ is the explicit constant
  \begin{equation}\label{eq:K_constant}
    K = \frac{2\delta^4 p^2(1-p)^2}{\sigma^6\gamma^2(1-\lambda)^2}\,
    (4\sigma^2 + \mu_2^2).
  \end{equation}
\end{enumerate}
\end{theorem}

\begin{corollary}[Asymmetric amplification]
\label{cor:asymmetric}
The bad-state amplification is exponentially larger than the good-state
dampening.
Specifically:
\begin{enumerate}[label=(\roman*)]
\item The good-state ruin ratio satisfies
  $\Prob_{\mu_1}(\tau^{\PI} \le T)\big/
   \Prob_{\mu_1}(\tau_1^{\FI} \le T) \to 0$
  as $a \to \infty$.
\item The bad-state ruin ratio satisfies, for $\epsilon < \delta p/\mu_2$:
  \begin{equation}\label{eq:bad_ratio_bound}
    \frac{\Prob_{\mu_2}(\tau^{\PI} < \infty)}
         {\Prob_{\mu_2}(\tau_2^{\FI} < \infty)}
    \;\ge\;
    C\,\exp\!\biggl(2\gamma(1-\lambda)\,a\,
      \frac{\delta p}{\delta p + \mu_2}\biggr)
  \end{equation}
  for an explicit constant $C \in (0,1)$.
  The amplification grows exponentially in the safety margin~$a$.
\end{enumerate}
At the \citeauthor{HuangSun2025} parameters
($\mu_1 = 0.2$, $\mu_2 = 0.02$, $\sigma = 0.1$, $\gamma = 5$,
$x = 3$, $D = 2$, $p = 0.5$, $\lambda = 0$):
\begin{align*}
  \text{Good state:} \quad
  & \Prob_{\mu_1}(\tau^{\PI})/\Prob_{\mu_1}(\tau_1^{\FI})
    \approx 0 \quad (\text{PI essentially ruin-free}), \\
  \text{Bad state:} \quad
  & \Prob_{\mu_2}(\tau^{\PI} < \infty)/\Prob_{\mu_2}(\tau_2^{\FI} < \infty)
    \approx 307.
\end{align*}
\end{corollary}

\begin{proposition}[Critical prior]\label{prop:critical_prior}
Define the mixture measure $\Prob_p := p\,\Prob_{\mu_1} + (1-p)\,\Prob_{\mu_2}$
and set
\[
  A := \Prob_{\mu_1}(\tau^{\PI} \le T) - \Prob_{\mu_1}(\tau_1^{\FI} \le T) \le 0
  \quad\text{(PI safer under good state)},
\]
\[
  B := \Prob_{\mu_2}(\tau^{\PI} \le T) - \Prob_{\mu_2}(\tau_2^{\FI} \le T) \ge 0
  \quad\text{(PI riskier under bad state)}.
\]
Then the critical prior at which the mixture comparison reverses is
\begin{equation}\label{eq:critical_prior}
  p^* := \frac{B}{B + |A|} \;\in\; (0,1).
\end{equation}
\textup{(i)} If $p > p^*$: $\Prob_p(\tau^{\PI} \le T) < \Prob_p(\tau^{\FI} \le T)$
  \textup{(PI safer under the mixture)}.
\textup{(ii)} If $p < p^*$: the inequality reverses \textup{(PI riskier)}.
The Lundberg approximation gives
$p^* \approx 1/\bigl(1 + e^{(\psi^{\PI} - \psi^{\FI})\,a}\bigr) \approx 0$
since $\psi^{\PI} \approx 5$--$6\times\psi^{\FI}$,
confirming that PI is safer under the mixture for essentially any non-degenerate prior.
Monte Carlo yields $p^* < 0.07$ across all $(D,\lambda)$ tested~\textup{(}Table~\ref{tab:bad_state}\textup{)}.
\end{proposition}

\begin{proof}
Set $A := \Prob_{\mu_1}(\tau^{\PI}) - \Prob_{\mu_1}(\tau^{\FI}) \le 0$ and
$B := \Prob_{\mu_2}(\tau^{\PI}) - \Prob_{\mu_2}(\tau^{\FI}) \ge 0$.
Then $\Prob_p(\tau^{\PI}) - \Prob_p(\tau^{\FI}) = pA + (1-p)B$,
which vanishes at $p^* = B/(B+|A|)$.
Since $|A| \gg B$ at realistic parameters (good-state dampening dominates
bad-state amplification), $p^*$ is small.
\end{proof}

\subsection{Sensitivity amplification}
\label{sec:sensitivity}

\begin{theorem}[$\lambda$-sensitivity comparison]
\label{thm:sensitivity}
The FI ruin probability under $\Prob_{\mu_k}$ satisfies
\begin{equation}\label{eq:FI_logP}
  -\log \Prob_{\mu_k}\bigl(\tau_k^{\FI} < \infty\bigr)
  = 2\gamma(1-\lambda)\,a,
\end{equation}
so that
$\frac{d}{d\lambda}\bigl[-\log \Prob_{\mu_k}(\tau_k^{\FI} < \infty)\bigr]
= -2\gamma\,a$.

The PI exponential bound (Theorem~\ref{thm:expo_bound}) gives
$-\log \Prob_{\mu_1}(\tau^{\PI} \le T) \ge 2\gamma(1-\lambda)\,a$,
with the effective PI Lundberg exponent satisfying the myopic lower bound
$\psi^{\PI} \geq 2\gamma(1-\lambda)\mu_1/\theta(p)$, giving a
sensitivity ratio $\psi^{\PI}/\psi^{\FI} \geq \mu_1/\theta(p)$.
Since $\pi^*(t) = g(t,P(t))/[\gamma(1-\lambda)]$ and $g$ is
$\lambda$-independent (Remark~\ref{rem:g_indep}), the PI portfolio
scales as $1/(1-\lambda)$, and the PI Lundberg exponent scales as
\begin{equation}\label{eq:PI_psi}
  \psi^{\PI}(\lambda) = \psi^{\PI}(0)\,(1-\lambda)\,,
  \qquad
  \frac{d\psi^{\PI}}{d\lambda} = -\psi^{\PI}(0).
\end{equation}
Because $\psi^{\PI}(0) \gg \psi^{\FI}(0) = 2\gamma$, the absolute
$\lambda$-sensitivity of the log-ruin probability is larger for PI.
More precisely, for the exponential ruin probability
$\Prob \approx e^{-\psi a}$:
\begin{equation}\label{eq:sensitivity_ratio}
  \frac{d/d\lambda\bigl[-\log \Prob(\tau^{\PI} \le T)\bigr]}
       {d/d\lambda\bigl[-\log \Prob(\tau_k^{\FI} < \infty)\bigr]}
  \;\approx\;
  \frac{\psi^{\PI}(0)}{\psi^{\FI}(0)}
  \;\approx\; 5\text{--}6.
\end{equation}
\end{theorem}

\begin{remark}\label{rem:sensitivity_paradox}
Theorem~\ref{thm:sensitivity} reveals a subtle paradox:
the PI investor is absolutely safer (lower ruin probability) but
proportionally more vulnerable to competition.
Increasing $\lambda$ from $0$ to $0.5$ multiplies the FI ruin
probability by $e^{2\gamma \cdot 0.5 \cdot a}$ but multiplies the PI
ruin probability by approximately
$e^{\psi^{\PI}(0) \cdot 0.5 \cdot a}$, a much larger factor.
The PI investor's ``buffer'' against ruin, while initially very large,
erodes faster with competition.
This is the precise mechanism behind the co-drawdown amplification
observed in \citet[Figure~4]{HuangSun2025}.
\end{remark}

\subsection{Heterogeneous investors}
\label{sec:heterogeneous}

We extend the analysis to a two-type game with heterogeneous risk
aversion.

\begin{proposition}[Heterogeneous investor perturbation]
\label{prop:heterogeneous}
Consider the two-type game with $N_A$ investors of risk aversion
$\gamma_A = \gamma + \varepsilon$ and $N_B$ investors of risk aversion
$\gamma_B = \gamma - \varepsilon$, where $\varepsilon \in (-\gamma, \gamma)$
is the heterogeneity parameter.
Let $\alpha_A = N_A/N$ and $\alpha_B = N_B/N$ be the population
fractions.
Then:

\begin{enumerate}[label=(\roman*)]
\item ($\lambda$-dependence at first order.)
  The PI strategy function $g_k(t,q;\lambda)$ for type~$k$ satisfies
  \begin{equation}\label{eq:g_hetero}
    g_k(t,q;\lambda) = g_0(t,q) + \varepsilon\,g_1^k(t,q;\lambda)
      + O(\varepsilon^2),
  \end{equation}
  where $g_0$ is the symmetric (and $\lambda$-independent) function
  from~\eqref{eq:g_decomp}, and $g_1^k$ depends on $\lambda$ through
  the linearized Cauchy problem
  \begin{equation}\label{eq:cauchy_linearized}
    \partial_t c_1^k + \frac{\beta^2}{2}\,\partial_{qq} c_1^k
    - \frac{\beta}{\sigma}\bigl[\theta - \sigma^2 \bar{g}_0\bigr]
      \partial_q c_1^k
    + \beta\sigma\,\bar{g}_1\,\partial_q c_0
    + \Delta\kappa_k\,\frac{\theta^2}{\sigma^2} = 0,
  \end{equation}
  with $c_1^k(T,q) = 0$,
  $\bar{g}_0 = g_0/[\gamma(1-\lambda)]$,
  $\Delta\kappa_A = -(1-\lambda/N)/\gamma^2$, and $\bar{g}_1$ is the
  first-order correction to the aggregate function depending on
  $g_1^A$, $g_1^B$, and $\lambda$.

\item (Wealth gap correction.)
  The expected wealth gap for type~$A$ satisfies
  \begin{equation}\label{eq:gap_hetero}
    \frac{\partial G_A}{\partial\lambda}
    = \frac{G_A(\lambda, T)}{1-\lambda}
    + \varepsilon\,\Phi(\lambda, T) + O(\varepsilon^2),
  \end{equation}
  where $\Phi(\lambda, T)$ captures the heterogeneity correction and
  is determined by the first-order Cauchy problem~\eqref{eq:cauchy_linearized}.

\item (Decorrelation effect.)
  Under PI, heterogeneity $\varepsilon \neq 0$ implies
  $\pi_A^*(t) \neq \pi_B^*(t)$ at first order.
  The wealth correlation between types satisfies
  $\rho_{AB} = 1 - O(\varepsilon^2)$, so the co-drawdown probability
  for a mixed population is strictly less than for the symmetric case.
\end{enumerate}
\end{proposition}

\begin{remark}\label{rem:hetero_mechanism}
The key mechanism is that heterogeneity breaks the cancellation in
Remark~\ref{rem:g_indep}.
When $\gamma_A \neq \gamma_B$, the aggregate strategy $\bar{g}$
involves a $\gamma_k$-weighted average where the more risk-tolerant
type contributes disproportionately.
This weighted average does not simplify to $g_0/[\gamma(1-\lambda)]$,
so $\lambda$ enters the Cauchy problem through the aggregate drift.
The resulting $\lambda$-dependence of $g_1^k$ is the first-order
symmetry-breaking effect of heterogeneity.
\end{remark}

\begin{remark}[Role of the symmetry assumption]\label{rem:symmetry_role}
The factorization $G(\lambda,T) = \mu_1 I(T)/[\gamma(1-\lambda)]$
and the $\lambda$-independence of $g(t,q)$ both require identical
$(\gamma, \lambda, x)$ across investors.
Proposition~\ref{prop:heterogeneous} shows that heterogeneity in
$\gamma$ alone introduces $\lambda$-dependence in~$g$ at first order.
The exponential safety bound (Theorem~\ref{thm:expo_bound}) does NOT
require symmetry: it holds for each individual investor.
\end{remark}

\begin{remark}[Role of constant drift]\label{rem:constant_drift}
The binary prior $\mu \in \{\mu_1, \mu_2\}$ with constant drift is
essential for the filter $P(t)$ to be a diffusion on~$[0,1]$.
Under regime-switching (alternating drift), the filter becomes
two-dimensional, and the hedging demand structure changes qualitatively.
The core mechanism (smaller $\pi$ $\Rightarrow$ higher
$\mu/(\pi\sigma^2)$ $\Rightarrow$ lower ruin) persists.
\end{remark}

\begin{remark}[Role of $r = 0$]\label{rem:zero_rate}
With $r > 0$, the equilibrium strategies acquire time-dependence
through $e^{-r(T-t)}$, breaking the stationarity used in the
Lundberg exponent computation.
The Dubins--Schwarz comparison (Theorem~\ref{thm:Q_comparison})
extends without modification since it depends only on the quadratic
variation ordering, not on stationarity.
\end{remark}

%% ========================================================================
\section{Proofs}\label{sec:proofs}
%% ========================================================================

\subsection{Proof of Theorem~\ref{thm:expo_bound}
  (Exponential safety bound)}
\label{sec:proof_expo}

The proof constructs an exponential supermartingale with a carefully
chosen tilting parameter.

\begin{proof}
Fix $T > 0$.
Under $\Prob_{\mu_1}$, the PI wealth process satisfies
\begin{equation}\label{eq:XPI_drift}
  X^{\PI}(t) = x + \int_0^t \pi^*(s)\,\mu_1\,ds
             + \int_0^t \pi^*(s)\,\sigma\,dW(s),
\end{equation}
where $\pi^*(s) = \theta(P(s))/[\sigma^2\gamma(1-\lambda)]$.
Since $\theta(q) \le \mu_1$ for $q \in [0,1]$, we have
\begin{equation}\label{eq:pi_upper}
  \pi^*(s) \le \frac{\mu_1}{\sigma^2\gamma(1-\lambda)}
  = \pi_1^{\FI} =: \pi_{\max}
  \qquad \text{a.s., for all } s \ge 0.
\end{equation}

Define the tilting parameter
\begin{equation}\label{eq:theta_tilt}
  \vartheta := -\frac{2\mu_1}{\pi_{\max}\,\sigma^2}
\end{equation}
and the exponential process
\begin{equation}\label{eq:Z_process}
  Z(t) := \exp\!\biggl(
    \vartheta\bigl(X^{\PI}(t) - x\bigr)
    - \int_0^t \psi(s)\,ds
  \biggr),
\end{equation}
where
\begin{equation}\label{eq:psi_def}
  \psi(s) := \vartheta\,\pi^*(s)\,\mu_1
           + \frac{\vartheta^2}{2}\,\bigl(\pi^*(s)\bigr)^2\,\sigma^2.
\end{equation}

\medskip
\noindent\textbf{Step~1: $Z$ is a local martingale.}
By It\^o's formula applied to $f(y) = e^{\vartheta y}$ with
$y(t) = X^{\PI}(t) - x - \int_0^t \psi(s)/\vartheta\,ds$:
\begin{align}
  dZ(t)
  &= Z(t)\,\vartheta\,\pi^*(t)\,\sigma\,dW(t). \label{eq:dZ}
\end{align}
More precisely, we verify this by writing
$Z(t) = \exp(\vartheta\int_0^t \pi^*(s)[\mu_1\,ds + \sigma\,dW(s)]
  - \int_0^t \psi(s)\,ds)$.
By the It\^o exponential:
\begin{align*}
  dZ(t)
  &= Z(t)\Bigl[\vartheta\,\pi^*(t)\,\mu_1\,dt
    + \vartheta\,\pi^*(t)\,\sigma\,dW(t)
    + \frac{\vartheta^2}{2}\,(\pi^*(t))^2\,\sigma^2\,dt
    - \psi(t)\,dt\Bigr] \\
  &= Z(t)\,\vartheta\,\pi^*(t)\,\sigma\,dW(t),
\end{align*}
since the $dt$ terms cancel by definition of $\psi$.
Therefore $Z$ is a non-negative local martingale, hence a
supermartingale with $\E[Z(t)] \le Z(0) = 1$.

\medskip
\noindent\textbf{Step~2: The compensator $\psi$ is non-positive.}
We show $\psi(s) \le 0$ for all $s$, which implies
$\int_0^t \psi(s)\,ds \le 0$ and therefore
$\exp(-\int_0^t \psi(s)\,ds) \ge 1$.

Factor $\psi(s)$:
\begin{align}
  \psi(s)
  &= \pi^*(s)\biggl[
    \vartheta\,\mu_1
    + \frac{\vartheta^2}{2}\,\pi^*(s)\,\sigma^2
  \biggr] \notag \\
  &= \pi^*(s)\,\frac{-2\mu_1^2}{\pi_{\max}\,\sigma^2}
    \biggl[1 - \frac{\pi^*(s)}{\pi_{\max}}\biggr],
  \label{eq:psi_factored}
\end{align}
where we used $\vartheta = -2\mu_1/(\pi_{\max}\sigma^2)$ and
$\vartheta^2\sigma^2/2 = 2\mu_1^2/(\pi_{\max}^2\sigma^2)$.

Since $\pi^*(s) \ge 0$ and $\pi^*(s) \le \pi_{\max}$ by~\eqref{eq:pi_upper},
the factor $[1 - \pi^*(s)/\pi_{\max}] \ge 0$, and the leading constant
$-2\mu_1^2/(\pi_{\max}\sigma^2) < 0$.
Therefore $\psi(s) \le 0$ for all $s \ge 0$.

\medskip
\noindent\textbf{Step~3: Optional stopping.}
Define $\tau := \tau^{\PI}$ for brevity.
Since $Z$ is a non-negative supermartingale:
\begin{equation}\label{eq:opt_stop}
  1 = Z(0) \ge \E_{\mu_1}\bigl[Z(\tau \wedge T)\bigr].
\end{equation}
On the event $\{\tau \le T\}$:
$X^{\PI}(\tau) = D$, so $X^{\PI}(\tau) - x = -(x-D) = -a$, and
\begin{equation}\label{eq:Z_at_tau}
  Z(\tau)
  = \exp\!\bigl(-\vartheta\,a - \textstyle\int_0^{\tau} \psi(s)\,ds\bigr)
  \ge \exp(-\vartheta\,a)
  = \exp\!\biggl(\frac{2\mu_1\,a}{\pi_{\max}\,\sigma^2}\biggr),
\end{equation}
where the inequality uses $\int_0^{\tau} \psi(s)\,ds \le 0$.

From~\eqref{eq:opt_stop}:
\begin{align}
  1
  &\ge \E_{\mu_1}\bigl[Z(\tau)\,\mathbf{1}_{\{\tau \le T\}}\bigr]
     + \E_{\mu_1}\bigl[Z(T)\,\mathbf{1}_{\{\tau > T\}}\bigr]
  \notag \\
  &\ge \exp\!\biggl(\frac{2\mu_1\,a}{\pi_{\max}\,\sigma^2}\biggr)
    \cdot \Prob_{\mu_1}(\tau \le T)
    + 0,
  \label{eq:opt_stop_bound}
\end{align}
using $Z(T) \ge 0$ and~\eqref{eq:Z_at_tau}.
Rearranging:
\begin{equation}\label{eq:main_bound}
  \Prob_{\mu_1}\bigl(\tau^{\PI} \le T\bigr)
  \le \exp\!\biggl(-\frac{2\mu_1\,a}{\pi_{\max}\,\sigma^2}\biggr).
\end{equation}

Finally, the classical first passage result for the FI process
$X_1^{\FI}(t) = x + \pi_1^{\FI}\mu_1 t + \pi_1^{\FI}\sigma W(t)$
(Brownian motion with positive drift $b = \pi_1^{\FI}\mu_1$ and
diffusion $c = \pi_1^{\FI}\sigma$) gives
\begin{equation}\label{eq:FI_infinite}
  \Prob_{\mu_1}\bigl(\tau_1^{\FI} < \infty\bigr)
  = \exp\!\biggl(-\frac{2b\,a}{c^2}\biggr)
  = \exp\!\biggl(-\frac{2\mu_1\,a}{\pi_1^{\FI}\,\sigma^2}\biggr).
\end{equation}
Since $\pi_{\max} = \pi_1^{\FI}$, \eqref{eq:main_bound}
and~\eqref{eq:FI_infinite} yield~\eqref{eq:expo_bound}.
\end{proof}

\subsection{Proof of Theorem~\ref{thm:Q_comparison}
  (Martingale measure comparison)}
\label{sec:proof_Q}

\begin{proof}
\textbf{Step~1: Girsanov transformation.}
Define $Q$ by~\eqref{eq:Q_density}.
Under $Q$, $\widetilde{W}(t) := W(t) + (\mu_1/\sigma)\,t$ is a
standard Brownian motion.
The stock has zero drift under $Q$:
$dS/S = \sigma\,d\widetilde{W}$.

The wealth processes become:
\begin{align}
  X_1^{\FI}(t) - x
    &= \pi_1^{\FI}\,\sigma\,\widetilde{W}(t),
  \label{eq:XFI_Q} \\
  X^{\PI}(t) - x
    &= \sigma\int_0^t \pi^*(s)\,d\widetilde{W}(s).
  \label{eq:XPI_Q}
\end{align}
Both are continuous $Q$-local martingales starting at~$0$.

\medskip
\noindent\textbf{Step~2: Dubins--Schwarz representation.}
By the Dambis--Dubins--Schwarz theorem
(\citealt[Theorem~3.4.6]{KaratzasShreve1991}), there exist standard
$Q$-Brownian motions $B_1$ and $B_2$ (possibly defined on an
extension of the probability space) such that:
\begin{align}
  X_1^{\FI}(t) - x &= B_1\bigl(U(t)\bigr),
    \qquad U(t) := \bigl(\pi_1^{\FI}\,\sigma\bigr)^2\,t,
  \label{eq:DS_FI} \\
  X^{\PI}(t) - x &= B_2\bigl(V(t)\bigr),
    \qquad V(t) := \sigma^2\int_0^t \bigl[\pi^*(s)\bigr]^2\,ds.
  \label{eq:DS_PI}
\end{align}
The clock $U$ is deterministic and linear.
The clock $V$ is random and satisfies
\begin{equation}\label{eq:V_le_U}
  V(T) = \sigma^2\int_0^T \bigl[\pi^*(s)\bigr]^2\,ds
  \le \sigma^2\,\bigl(\pi_1^{\FI}\bigr)^2\,T = U(T),
  \qquad Q\text{-a.s.},
\end{equation}
since $\pi^*(s) \le \pi_1^{\FI}$ a.s.\ for all $s$.
The inequality is strict with positive $Q$-probability whenever
$\pi^*(s) < \pi_1^{\FI}$ on a set of positive Lebesgue measure, which
holds for $p < 1$.

\medskip
\noindent\textbf{Step~3: First passage time comparison.}
The ruin events translate to:
\begin{align}
  \bigl\{\tau_1^{\FI} \le T\bigr\}
    &= \Bigl\{\inf_{0 \le t \le T}\bigl(X_1^{\FI}(t) - x\bigr) \le -a\Bigr\}
     = \Bigl\{\inf_{0 \le u \le U(T)} B_1(u) \le -a\Bigr\},
  \label{eq:ruin_FI_DS} \\
  \bigl\{\tau^{\PI} \le T\bigr\}
    &= \Bigl\{\inf_{0 \le t \le T}\bigl(X^{\PI}(t) - x\bigr) \le -a\Bigr\}
     = \Bigl\{\inf_{0 \le v \le V(T)} B_2(v) \le -a\Bigr\}.
  \label{eq:ruin_PI_DS}
\end{align}

Since $V(T) \le U(T)$ a.s.:
\begin{equation}\label{eq:subset}
  \Bigl\{\inf_{0 \le v \le V(T)} B_2(v) \le -a\Bigr\}
  \;\subseteq\;
  \Bigl\{\inf_{0 \le v \le U(T)} B_2(v) \le -a\Bigr\}.
\end{equation}

Taking $Q$-probabilities and using the fact that $B_2$ under $Q$ has
the same law as $B_1$ under $Q$ (both are standard Brownian motions):
\begin{align}
  Q\bigl(\tau^{\PI} \le T\bigr)
  &\le Q\Bigl(\inf_{0 \le v \le U(T)} B_2(v) \le -a\Bigr) \notag \\
  &= Q\Bigl(\inf_{0 \le u \le U(T)} B_1(u) \le -a\Bigr) \notag \\
  &= Q\bigl(\tau_1^{\FI} \le T\bigr).
  \label{eq:Q_comparison_chain}
\end{align}

\medskip
\noindent\textbf{Step~4: Equality condition.}
Equality in~\eqref{eq:subset} holds if and only if $V(T) = U(T)$
$Q$-a.s., which by~\eqref{eq:V_le_U} requires
$\pi^*(s) = \pi_1^{\FI}$ for Lebesgue-a.e.\ $s \in [0,T]$, $Q$-a.s.
\end{proof}

\begin{remark}\label{rem:DS_subtlety}
In Step~3, we used the inclusion~\eqref{eq:subset} rather than a
direct coupling of $B_1$ and $B_2$.
This is important because the Dubins--Schwarz Brownian motions $B_1$
and $B_2$ are constructed from the respective martingales and are
\emph{not} identical processes in general.
However, their marginal laws under $Q$ are identical (both are standard
Brownian motions), which suffices for the probability comparison.
\end{remark}

\begin{remark}[Quantifying the Girsanov gap]\label{rem:girsanov_gap}
The gap between Theorem~\ref{thm:Q_comparison} (under~$Q$) and
Conjecture~\ref{conj:physical} (under $\Prob_{\mu_1}$) satisfies
\[
  \bigl|\Prob_{\mu_1}(\tau^I \le T) - Q(\tau^I \le T)\bigr|
  \le \sqrt{e^{m^2 T} - 1}\;\sqrt{Q(\tau^I \le T)}\,,
  \qquad m := \mu_1/\sigma,
\]
by Cauchy--Schwarz.
At the calibrated parameters ($m^2 T = 40$), this bound exceeds~$10^8$
and is thus useless.
A tighter analysis decomposes
$\Prob_{\mu_1}(\tau^I \le T)
  = \E_Q[L(T) \mid \tau^I \le T]\;Q(\tau^I \le T)$,
where $L(T) = e^{m\widetilde{W}(T) - m^2 T/2}$.
The conjecture reduces to verifying
$\E_Q[L(T) \mid \tau^{\PI} \le T] \le \E_Q[L(T) \mid \tau^{\FI} \le T]$:
the Girsanov weight, conditional on ruin, should be smaller for PI
because the PI ruin event requires a larger negative stock excursion
(smaller portfolio $\Rightarrow$ larger decline needed to reach the
barrier), and such paths carry smaller~$L(T)$.
\end{remark}

\subsection{Proof of Corollary~\ref{cor:constant}
  (Constant damping)}
\label{sec:proof_constant}

\begin{proof}
When $\pi^*(t) \equiv c$ for a constant $c \in (0, \pi_1^{\FI})$, the
PI wealth process under any measure~$\Prob$ equivalent to $\Prob_{\mu_1}$
satisfies
\begin{equation}\label{eq:constant_wealth}
  X^{\PI}(t) - x = c\bigl(\mu_1\,t + \sigma\,W(t)\bigr)
  = \frac{c}{\pi_1^{\FI}}\,\bigl(X_1^{\FI}(t) - x\bigr).
\end{equation}
Define $\alpha := c/\pi_1^{\FI} \in (0,1)$.
Then $X^{\PI}(t) = (1-\alpha)\,x + \alpha\,X_1^{\FI}(t)$ and
\begin{equation}\label{eq:constant_ruin}
  X^{\PI}(t) \le D
  \;\;\Longleftrightarrow\;\;
  X_1^{\FI}(t) \le \frac{D - (1-\alpha)\,x}{\alpha}
  = x - \frac{a}{\alpha}\,.
\end{equation}
Since $\alpha < 1$, the effective barrier $x - a/\alpha < x - a = D$ is
strictly below~$D$.
Therefore
\begin{equation}\label{eq:constant_inclusion}
  \bigl\{\tau^{\PI} \le T\bigr\}
  = \Bigl\{\inf_{t \le T} X_1^{\FI}(t) \le x - \frac{a}{\alpha}\Bigr\}
  \;\subsetneq\;
  \bigl\{\tau_1^{\FI} \le T\bigr\},
\end{equation}
and the inclusion is strict because the event
$\{x - a/\alpha < \inf_{t \le T} X_1^{\FI}(t) \le D\}$ has positive
probability.
Taking probabilities under any measure gives~\eqref{eq:const_comp}.
\end{proof}

%% ========================================================================
%% Section 4 (continued): Proofs
%% ========================================================================

\subsection{Proof of Proposition~\ref{prop:over_invest}}\label{sec:proof_over_invest}

We prove that under $\Prob_{\mu_2}$ (the bad state), the PI investor
over-invests relative to the FI optimal: $\E_{\mu_2}[\pi^*(t)] > \pi_2^{\FI}$
for all $t > 0$ and $p \in (0,1)$.

\begin{proof}
The proof proceeds in four steps.

\medskip\noindent\textbf{Step 1: Filter dynamics under $\Prob_{\mu_2}$.}
Under $\Prob_{\mu_2}$, the stock return is $dR(t) = \mu_2\,dt + \sigma\,dW(t)$,
where $W$ is a standard $\Prob_{\mu_2}$-Brownian motion.
The Bayesian filter~\eqref{eq:filter} admits the Bayes formula
representation
\begin{equation}\label{eq:filter_bad}
  P(t) = \frac{1}{1 + \frac{1-p}{p}\,\Lambda(t)},
\end{equation}
where $\Lambda(t) := \ell_2(t)/\ell_1(t)$ is the likelihood ratio with
\begin{equation}\label{eq:Lambda_bad}
  \Lambda(t) = \exp\!\biggl(-\frac{\delta\,W(t)}{\sigma}
  + \frac{\delta^2\,t}{2\sigma^2}\biggr), \qquad \delta := \mu_1 - \mu_2.
\end{equation}
Indeed, the likelihood functions satisfy
$\ell_k(t) = \exp\bigl(\mu_k R(t)/\sigma^2 - \mu_k^2 t/(2\sigma^2)\bigr)$.
Substituting $R(t) = \mu_2 t + \sigma W(t)$ and simplifying:
\[
  \frac{\ell_2(t)}{\ell_1(t)}
  = \exp\!\biggl(\frac{(\mu_2 - \mu_1)R(t)}{\sigma^2}
  + \frac{\mu_1^2 - \mu_2^2}{2\sigma^2}\,t\biggr)
  = \exp\!\biggl(-\frac{\delta\,W(t)}{\sigma}
  + \frac{\delta^2\,t}{2\sigma^2}\biggr),
\]
using $(\mu_1+\mu_2)/2 - \mu_2 = \delta/2$.

\medskip\noindent\textbf{Step 2: $P(t) > 0$ almost surely for finite $t$.}
Since $W(t)$ is finite a.s.\ for each fixed $t$, we have $\Lambda(t) < \infty$
a.s., hence $P(t) = 1/(1 + \frac{1-p}{p}\Lambda(t)) > 0$ a.s.\ for all
$t \in [0,\infty)$.
As $t \to \infty$, the strong law of large numbers gives $W(t)/t \to 0$ a.s.,
so $\Lambda(t) \to +\infty$ a.s.\ (the drift $\delta^2 t/(2\sigma^2)$
dominates), yielding $P(t) \to 0$ a.s.
However, this convergence requires infinite time; for any finite~$t$,
$P(t) > 0$ a.s.

\medskip\noindent\textbf{Step 3: The filter is a $\Prob_{\mu_2}$-supermartingale
with $\E_{\mu_2}[P(t)] > 0$.}
Under the mixture measure $\Prob_p := p\,\Prob_{\mu_1} + (1-p)\,\Prob_{\mu_2}$,
the filter $P(t) = \Prob_p(\mu = \mu_1 \mid \mathcal{F}_t^S)$ is a
$\Prob_p$-martingale by the tower property:
\begin{equation}\label{eq:filter_martingale}
  p = \E_{\Prob_p}[P(t)]
  = p\,\E_{\mu_1}[P(t)] + (1-p)\,\E_{\mu_2}[P(t)].
\end{equation}
Under $\Prob_{\mu_1}$, the filter learns the truth:
$P(t) \to 1$ a.s.\ as $t \to \infty$, so $\E_{\mu_1}[P(t)] > p$ for $t > 0$
(the filter is a $\Prob_{\mu_1}$-submartingale, as established in
the Bayes formula~\eqref{eq:filter_bad}).
Substituting into~\eqref{eq:filter_martingale}:
$(1-p)\,\E_{\mu_2}[P(t)] < p - p\cdot p = p(1-p)$,
which gives $\E_{\mu_2}[P(t)] < p$.
Together with the pointwise bound $P(t) > 0$ a.s., we conclude
$\E_{\mu_2}[P(t)] \in (0, p)$ for all $t > 0$.

\medskip\noindent\textbf{Step 4: Over-investment.}
Since $\theta(q) = \delta q + \mu_2$ is affine:
\begin{equation}\label{eq:Etheta_bad}
  \E_{\mu_2}[\theta(P(t))] = \delta\,\E_{\mu_2}[P(t)] + \mu_2.
\end{equation}
By Step~3, $\E_{\mu_2}[P(t)] > 0$, hence
\[
  \E_{\mu_2}[\pi^*(t)]
  = \frac{\E_{\mu_2}[\theta(P(t))]}{\sigma^2\gamma(1-\lambda)}
  = \frac{\delta\,\E_{\mu_2}[P(t)] + \mu_2}{\sigma^2\gamma(1-\lambda)}
  > \frac{\mu_2}{\sigma^2\gamma(1-\lambda)}
  = \pi_2^{\FI}
\]
for all $t > 0$ and $p \in (0,1)$.
\end{proof}

\begin{remark}[Initial over-investment ratio]\label{rem:overinvest_quant}
At $t = 0$: $\pi^*(0)/\pi_2^{\FI} = (\delta p + \mu_2)/\mu_2 = 1 + \delta p/\mu_2$.
For the Huang--Sun parameters $(\delta = 0.18, p = 0.5, \mu_2 = 0.02)$,
this ratio equals~$5.5$.
The PI investor initially invests more than five times the FI optimal
under the bad state, because her prior assigns probability~$p = 0.5$ to the
good state $\mu_1 = 0.2$ while the true drift is only $\mu_2 = 0.02$.
\end{remark}


%% ------------------------------------------------------------------------

\begin{lemma}[Filter modulus of continuity]\label{lem:filter_modulus}
Under either $\Prob_{\mu_k}$, the filter satisfies
\[
  \E\bigl[|P(s) - P(t)|^2 \,\big|\, \mathcal{F}_t\bigr]
  \;\le\; \frac{\delta^2}{16\sigma^2}\,(s-t),
  \qquad 0 \le t \le s,
\]
and consequently
$\E\bigl[\sup_{s \in [t,t+h]} |P(s) - P(t)|^2\bigr]
  \le \delta^2 h/(4\sigma^2)$.
\end{lemma}

\begin{proof}
By It\^{o} isometry, $\E[|P(s)-P(t)|^2 \mid \mathcal{F}_t]
  = \E[\int_t^s \beta(P(u))^2\,du \mid \mathcal{F}_t]$
where $\beta(q) = (\delta/\sigma)q(1-q)$.
Since $\max_{q \in [0,1]} q^2(1-q)^2 = 1/16$, we have
$\beta(q)^2 \le \delta^2/(16\sigma^2)$.
The supremum bound follows from Doob's $L^2$-maximal inequality.
\end{proof}

\begin{remark}\label{rem:filter_modulus_reusable}
Lemma~\ref{lem:filter_modulus} applies to any binary-state Bayesian
filtering problem with Gaussian observations.
The bound $\delta^2/(16\sigma^2)$ depends only on the signal-to-noise
ratio $\delta/\sigma$ and is independent of the true state, prior, or
time.
This makes it a reusable tool for $L^2$-coupling arguments involving
Bayesian filters.
\end{remark}

\subsection{Proof of Theorem~\ref{thm:bad_state}}\label{sec:proof_bad_state_detail}
%% ------------------------------------------------------------------------

We prove the bad-state ruin comparison in three stages:
(A)~exact FI ruin probability,
(B)~PI ruin lower bound via $L^2$ coupling with a constant-portfolio reference,
and (C)~amplification ratio computation.

\begin{proof}
\textbf{Stage A: Exact FI ruin probability.}

Under $\Prob_{\mu_2}$, the FI wealth deviation
$Y_2^{\FI}(t) := X_2^{\FI}(t) - x = \pi_2^{\FI}(\mu_2 t + \sigma W(t))$
is Brownian motion with drift $b := \pi_2^{\FI}\mu_2 > 0$ and diffusion
$c := \pi_2^{\FI}\sigma > 0$.
By the classical first-passage formula
\citep[Proposition~3.5.2]{KaratzasShreve1991}, the infinite-horizon
ruin probability is
\begin{equation}\label{eq:FI_ruin_exact}
  \Prob_{\mu_2}(\tau_2^{\FI} < \infty)
  = \exp\!\biggl(-\frac{2b\,a}{c^2}\biggr)
  = \exp\!\biggl(-\frac{2\mu_2\,a}{\pi_2^{\FI}\sigma^2}\biggr)
  = \exp\!\bigl(-2\gamma(1-\lambda)\,a\bigr),
\end{equation}
where $a := x - D > 0$ is the safety margin and we used
$\pi_2^{\FI} = \mu_2/[\sigma^2\gamma(1-\lambda)]$.

\medskip
\textbf{Stage B: PI ruin lower bound via coupling.}

The argument compares the PI wealth process to a constant-portfolio
reference process over a short time window, using $L^2$ proximity.
We proceed through four substeps.

\medskip\noindent\textbf{Step B1: Reference process and BL monotonicity.}
Define the constant-portfolio reference wealth:
\begin{equation}\label{eq:bar_pi}
  \bar{\pi} := \pi^*(0) = \frac{\theta(p)}{\sigma^2\gamma(1-\lambda)}
  = \frac{\delta p + \mu_2}{\sigma^2\gamma(1-\lambda)}.
\end{equation}
Note that $\bar{\pi}/\pi_2^{\FI} = (\delta p + \mu_2)/\mu_2
= 1 + \delta p/\mu_2 > 1$.
The reference process is
$\widetilde{Z}(t) := \bar{\pi}(\mu_2 t + \sigma W(t))$,
which is Brownian motion with drift $\bar{\pi}\mu_2$ and
diffusion $\bar{\pi}\sigma$.

\begin{lemma}[Bachelier--L\'evy monotonicity]\label{lem:BL_mono}
For a constant portfolio $\pi > 0$ and the process
$Z_\pi(t) := \pi(\mu_2 t + \sigma B(t))$, the ruin probability
$\mathfrak{p}(a, \pi\mu_2, \pi\sigma, T)$
is strictly increasing in~$\pi$ for each $a > 0$ and $T > 0$.
\end{lemma}

\begin{proof}[Proof of Lemma~\ref{lem:BL_mono}]
Write $\inf_{t \le T} Z_\pi(t) = \pi\,\inf_{t \le T}(\mu_2 t + \sigma B(t))
=: \pi\,m_T$, where $m_T < 0$ a.s.\
(since Brownian motion takes negative values near time~0).
Then
$\{Z_\pi < -a\} = \{\pi\,m_T < -a\} = \{m_T < -a/\pi\}$.
Since $a/\pi$ is strictly decreasing in~$\pi$ and $\Prob(m_T < -c)$ is
strictly decreasing in~$c > 0$ (as $m_T$ has a continuous distribution),
$\mathfrak{p}(a, \pi, T)$ is strictly increasing in~$\pi$.
\end{proof}

By Lemma~\ref{lem:BL_mono}, since $\bar{\pi} > \pi_2^{\FI}$:
\begin{equation}\label{eq:BL_comparison}
  \Prob_{\mu_2}(\widetilde{\tau} \le T)
  = \mathfrak{p}(a, \bar{\pi}, T)
  > \mathfrak{p}(a, \pi_2^{\FI}, T)
  = \Prob_{\mu_2}(\tau_2^{\FI} \le T)
\end{equation}
for all $T > 0$, where $\widetilde{\tau} := \inf\{t \ge 0 :
\widetilde{Z}(t) < -a\}$.
In the limit $T \to \infty$:
\begin{equation}\label{eq:R_barpi_proof}
  R(\bar{\pi}) := \Prob_{\mu_2}(\widetilde{\tau} < \infty)
  = \exp\!\biggl(-\frac{2\mu_2\,a}{\bar{\pi}\sigma^2}\biggr)
  > \exp\!\biggl(-\frac{2\mu_2\,a}{\pi_2^{\FI}\sigma^2}\biggr)
  = \Prob_{\mu_2}(\tau_2^{\FI} < \infty).
\end{equation}

\medskip\noindent\textbf{Step B2: Filter modulus of continuity.}
Under $\Prob_{\mu_2}$, the filter satisfies the SDE
\begin{equation}\label{eq:filter_SDE_bad}
  dP(t) = \frac{\delta\,P(t)(1-P(t))}{\sigma}\,d\widehat{W}(t),
  \qquad P(0) = p,
\end{equation}
where $\widehat{W}$ is the innovation process (a standard Brownian motion
under $\Prob_{\mu_2}$).
Define $\sigma_P := \delta\,p(1-p)/\sigma$ (the initial filter volatility).
For short times $s \le 1$:
\begin{equation}\label{eq:filter_L2_proof}
  \E_{\mu_2}\bigl[|P(s) - p|^2\bigr]
  = \sigma_P^2\,s + O(s^2)
  \le C_P\,s,
\end{equation}
where $C_P := 2\delta^2 p^2(1-p)^2/\sigma^2$.

Since $\pi^*(t) = (\delta\,P(t) + \mu_2)/[\sigma^2\gamma(1-\lambda)]$
(myopic approximation), the strategy modulus satisfies
\begin{equation}\label{eq:pi_L2_proof}
  \E_{\mu_2}\bigl[|\pi^*(s) - \bar{\pi}|^2\bigr]
  = \frac{\delta^2}{\sigma^4\gamma^2(1-\lambda)^2}\,
    \E_{\mu_2}\bigl[|P(s) - p|^2\bigr]
  \le C_\pi\,s,
\end{equation}
where $C_\pi := 2\delta^4 p^2(1-p)^2/[\sigma^6\gamma^2(1-\lambda)^2]$.

\medskip\noindent\textbf{Step B3: $L^2$ bound on the process difference.}
Let $Z^{\PI}(t) := X^{\PI}(t) - x$ be the PI wealth deviation and define
$\Delta(t) := Z^{\PI}(t) - \widetilde{Z}(t)$.
Then
\[
  \Delta(t) = \int_0^t (\pi^*(s) - \bar{\pi})
  \bigl(\mu_2\,ds + \sigma\,dW(s)\bigr)
\]
is a semimartingale.
By Doob's $L^2$-maximal inequality applied to the martingale part
$M_\Delta(t) := \int_0^t (\pi^*(s) - \bar{\pi})\sigma\,dW(s)$
together with the It\^o isometry, and by Cauchy--Schwarz for the drift part
$A_\Delta(t) := \int_0^t (\pi^*(s) - \bar{\pi})\mu_2\,ds$:
\begin{equation}\label{eq:Delta_L2}
  \E_{\mu_2}\!\Bigl[\sup_{s \le h}|\Delta(s)|^2\Bigr]
  \le 4\sigma^2 C_\pi\,h^2 + \mu_2^2 C_\pi\,h^3
  \le K\,h^2
\end{equation}
for $h \le 1$, where $K := C_\pi(4\sigma^2 + \mu_2^2)$.

By Markov's inequality:
\begin{equation}\label{eq:Delta_Markov}
  \Prob_{\mu_2}\!\Bigl(\sup_{s \le h}|\Delta(s)| > \eta\Bigr)
  \le \frac{K\,h^2}{\eta^2}.
\end{equation}

\medskip\noindent\textbf{Step B4: Ruin probability approximation.}
For any $\eta > 0$ and $h > 0$, we have the event inclusion
\[
  \bigl\{\inf_{s \le h}\widetilde{Z}(s) < -(a+\eta)\bigr\}
  \cap \bigl\{\sup_{s \le h}|\Delta(s)| \le \eta\bigr\}
  \;\subseteq\;
  \bigl\{\inf_{s \le h} Z^{\PI}(s) < -a\bigr\}.
\]
Indeed, on this event: $\inf Z^{\PI} \le \inf\widetilde{Z} + \sup|\Delta|
< -(a+\eta) + \eta = -a$.
Taking probabilities and using $\Prob(A \cap B) \ge \Prob(A) - \Prob(B^c)$:
\begin{equation}\label{eq:Main_LB}
  \Prob_{\mu_2}(\tau^{\PI} \le h)
  \ge \mathfrak{p}\!\bigl((1+\epsilon)a,\;\bar{\pi}\mu_2,\;
      \bar{\pi}\sigma,\;h\bigr)
  - \frac{K\,h^2}{\epsilon^2 a^2},
\end{equation}
where we set $\eta = \epsilon\,a$ for a parameter
$\epsilon \in (0, \delta p/\mu_2)$.

\medskip
\textbf{Stage C: Amplification ratio.}

We now show that the right-hand side of~\eqref{eq:Main_LB} exceeds
$\Prob_{\mu_2}(\tau_2^{\FI} < \infty)$ for appropriate~$h$.

By the Bachelier--L\'evy formula
\citep[Proposition~3.5.2]{KaratzasShreve1991}:
\begin{equation}\label{eq:BL_formula}
  \mathfrak{p}(a', b, c, T)
  = \Phi\!\biggl(\frac{-a'-bT}{c\sqrt{T}}\biggr)
  + e^{-2a'b/c^2}\,\Phi\!\biggl(\frac{-a'+bT}{c\sqrt{T}}\biggr),
\end{equation}
where $\Phi$ is the standard normal CDF.
As $T \to \infty$: the first $\Phi$-term converges to~$0$ (argument $\to -\infty$)
and the second $\Phi$-term converges to~$1$ (argument $\to +\infty$),
so $\mathfrak{p}(a', b, c, T) \to e^{-2a'b/c^2}$.

With $a' = (1+\epsilon)a$, $b = \bar{\pi}\mu_2$, $c = \bar{\pi}\sigma$:
\begin{equation}\label{eq:R_epsilon}
  R_\epsilon := \lim_{T \to \infty}\mathfrak{p}\!\bigl((1+\epsilon)a,\;
    \bar{\pi}\mu_2,\;\bar{\pi}\sigma,\;T\bigr)
  = \exp\!\biggl(-\frac{2(1+\epsilon)\mu_2\,a}{\bar{\pi}\sigma^2}\biggr).
\end{equation}

We need $R_\epsilon > R(\pi_2^{\FI})$, i.e.,
$(1+\epsilon)/\bar{\pi} < 1/\pi_2^{\FI}$, i.e.,
$\epsilon < \bar{\pi}/\pi_2^{\FI} - 1 = \delta p/\mu_2$.
This holds by hypothesis.
Define $\Delta R := R_\epsilon - R(\pi_2^{\FI}) > 0$.

\emph{Choosing~$h$:}
Since $\mathfrak{p}$ is continuous and converges to $R_\epsilon$ as
$T \to \infty$, there exists $h_1 > 0$ such that for all $h \ge h_1$:
\[
  \mathfrak{p}\!\bigl((1+\epsilon)a,\;\bar{\pi}\mu_2,\;
    \bar{\pi}\sigma,\;h\bigr)
  \ge R_\epsilon - \Delta R/4.
\]
Choose $h_2 > 0$ such that $K h_2^2/(\epsilon^2 a^2) = \Delta R/4$,
i.e., $h_2 = \epsilon\,a\,\sqrt{\Delta R/(4K)}$.
Setting $h = \max(h_1, h_2)$ in~\eqref{eq:Main_LB}:
\begin{align*}
  \Prob_{\mu_2}(\tau^{\PI} \le h)
  &\ge \bigl(R_\epsilon - \Delta R/4\bigr) - \Delta R/4 \\
  &= R_\epsilon - \Delta R/2 \\
  &> R(\pi_2^{\FI})
  = \Prob_{\mu_2}(\tau_2^{\FI} < \infty).
\end{align*}

Since $\Prob_{\mu_2}(\tau^{\PI} < \infty) \ge \Prob_{\mu_2}(\tau^{\PI} \le h)$,
this establishes part~(i) of the theorem.

\medskip\noindent\textbf{Amplification ratio.}
Taking $\epsilon \to 0$ in the ratio of ruin probabilities:
\begin{equation}\label{eq:amp_ratio_proof}
  \frac{\Prob_{\mu_2}(\tau^{\PI} < \infty)}
       {\Prob_{\mu_2}(\tau_2^{\FI} < \infty)}
  \ge C \cdot \frac{R(\bar{\pi})}{R(\pi_2^{\FI})}
  = C \cdot \exp\!\biggl(\frac{2\mu_2\,a}{\sigma^2}
    \biggl(\frac{1}{\pi_2^{\FI}} - \frac{1}{\bar{\pi}}\biggr)\biggr),
\end{equation}
where $C \in (0,1)$ accounts for the $L^2$~error terms.
The exponent simplifies:
\begin{align}\label{eq:ratio_exponent}
  \frac{2\mu_2\,a}{\sigma^2}\biggl(\frac{1}{\pi_2^{\FI}} - \frac{1}{\bar{\pi}}\biggr)
  &= \frac{2\mu_2\,a}{\sigma^2}\cdot\sigma^2\gamma(1-\lambda)
    \cdot\frac{\delta p}{\mu_2(\delta p + \mu_2)} \notag\\
  &= 2\gamma(1-\lambda)\,a\,\frac{\delta p}{\delta p + \mu_2}.
\end{align}
This completes the proof.
\end{proof}

\begin{remark}[Numerical illustration of the amplification ratio]
\label{rem:amp_numerical}
For the Huang--Sun parameters $(\mu_1 = 0.2, \mu_2 = 0.02, \sigma = 0.1,
\gamma = 5, \lambda = 0.3, p = 0.5, a = 1)$:
\begin{itemize}
\item $\bar{\pi}/\pi_2^{\FI} = 5.5$;
\item $R(\pi_2^{\FI}) = e^{-7.00} \approx 9.1 \times 10^{-4}$;
\item $R(\bar{\pi}) = e^{-1.27} \approx 0.280$;
\item Exponential amplification:
  $R(\bar{\pi})/R(\pi_2^{\FI}) = e^{5.73} \approx 307$.
\end{itemize}
Even with a conservative constant $C = 0.1$ in~\eqref{eq:amp_ratio_proof},
the PI ruin probability exceeds the FI one by a factor of at
least~$30$.
\end{remark}

\begin{remark}[Comparison with finite-horizon numerics]
\label{rem:finite_horizon}
The infinite-horizon bound in Theorem~\ref{thm:bad_state} yields an
amplification ratio of~$\sim 307$ at $a = 1$, while the Monte Carlo
simulations at finite horizon $T = 10$ (Table~\ref{tab:bad_state})
show ratios of $2$--$3$.
The discrepancy arises because the finite-horizon ruin probability is
substantially smaller than the infinite-horizon one, and the coupling error
term $K h^2/(\epsilon^2 a^2)$ is non-negligible for moderate~$h$.
The theorem correctly predicts the \emph{direction} (PI riskier under
the bad state) and the \emph{mechanism} (over-investment), while the
precise finite-horizon ratio requires numerical computation.
\end{remark}


%% ------------------------------------------------------------------------
\subsection{Proof of Theorem~\ref{thm:sensitivity}}\label{sec:proof_sensitivity}
%% ------------------------------------------------------------------------

We prove that the PI ruin probability is more sensitive to the competition
parameter~$\lambda$ than the FI ruin probability, with an explicit
amplification factor.

\begin{proof}
The key structural fact, established in Remark~\ref{rem:g_indep},
is that the function~$g(t,q)$ does not depend on~$\lambda$ under
the symmetric game.
Consequently, the PI strategy is
$\pi^*(t) = g(t,P(t))/[\gamma(1-\lambda)]$,
and the dependence on~$\lambda$ enters \emph{only} through the
explicit prefactor $1/[\gamma(1-\lambda)]$.

\medskip\noindent\textbf{Step 1: FI sensitivity.}
Under $\Prob_{\mu_1}$ (the good state), the FI wealth is a drifted
Brownian motion with $\pi^{\FI} = \mu_1/[\sigma^2\gamma(1-\lambda)]$.
The infinite-horizon ruin probability is
\[
  \Prob_{\mu_1}(\tau^{\FI} < \infty)
  = \exp\!\biggl(-\frac{2\mu_1\,a}{\pi^{\FI}\sigma^2}\biggr)
  = e^{-2\gamma(1-\lambda)a}.
\]
Define $\psi^{\FI} := 2\gamma(1-\lambda)$, so that
$-\log\Prob_{\mu_1}(\tau^{\FI} < \infty) = \psi^{\FI}\cdot a$.
The sensitivity is
\begin{equation}\label{eq:FI_sensitivity}
  \frac{d}{d\lambda}\bigl(-\log\Prob(\tau^{\FI} < \infty)\bigr)
  = -2\gamma\,a.
\end{equation}
This is negative: increasing $\lambda$ increases FI ruin probability,
as expected from the portfolio amplification
$\pi^{\FI} \propto 1/(1-\lambda)$.

\medskip\noindent\textbf{Step 2: PI Lundberg exponent.}
For the PI investor, the effective ``Lundberg exponent'' governing the
exponential decay of ruin probability with the safety margin~$a$ is
\begin{equation}\label{eq:psi_PI}
  \psi^{\PI} := -\frac{\log\Prob_{\mu_1}(\tau^{\PI} < \infty)}{a}.
\end{equation}
By Theorem~\ref{thm:expo_bound}, $\psi^{\PI} \ge \psi^{\FI}$
(PI is exponentially safer under the good state).

The exponential martingale bound (Theorem~\ref{thm:expo_bound})
gives
\[
  \Prob_{\mu_1}(\tau^{\PI} \le T)
  \le \exp\!\biggl(-\frac{2\mu_1\,a}{\pi_{\max}\,\sigma^2}\biggr),
\]
where $\pi_{\max} = \mu_1/[\sigma^2\gamma(1-\lambda)]$
(the maximum possible PI portfolio, attained when $P(t) = 1$).
Since $\pi_{\max} = \pi^{\FI}$, the bound recovers the FI
infinite-horizon probability.

For the myopic approximation (ignoring the hedging term), the initial
PI portfolio is $\bar{\pi} = \theta(p)/[\sigma^2\gamma(1-\lambda)]$.
The effective Lundberg exponent for the PI investor under the myopic
approximation satisfies
\begin{equation}\label{eq:psi_PI_myopic}
  \psi^{\PI}_{\mathrm{myopic}} \approx \frac{2\mu_1}{\bar{\pi}_0\,\sigma^2}
  = \frac{2\mu_1\gamma(1-\lambda)}{\theta(p)},
\end{equation}
where $\bar{\pi}_0 := \theta(p)/[\sigma^2\gamma(1-\lambda)]$ is the
initial PI portfolio.
Since $\theta(p) = \delta p + \mu_2 < \mu_1$ for $p < 1$, we have
$\psi^{\PI}_{\mathrm{myopic}} > \psi^{\FI}$.

\medskip\noindent\textbf{Step 3: Sensitivity ratio.}
Both Lundberg exponents are proportional to $(1-\lambda)$:
\[
  \psi^{\FI} = 2\gamma(1-\lambda), \qquad
  \psi^{\PI}_{\mathrm{myopic}} = \frac{2\mu_1\gamma(1-\lambda)}{\theta(p)}.
\]
The ratio is
\begin{equation}\label{eq:sensitivity_ratio_proof}
  \frac{\psi^{\PI}_{\mathrm{myopic}}}{\psi^{\FI}}
  = \frac{\mu_1}{\theta(p)}
  = \frac{\mu_1}{\delta p + \mu_2}.
\end{equation}
For the Huang--Sun parameters: $\mu_1/\theta(0.5) = 0.2/0.11 \approx 1.82$.

The \emph{sensitivity} of $-\log\Prob(\tau < \infty)$ to~$\lambda$ is
proportional to $\psi$, so the PI sensitivity exceeds FI by the same factor:
\begin{equation}\label{eq:sensitivity_amp}
  \frac{d(-\log\Prob(\tau^{\PI}))/d\lambda}
       {d(-\log\Prob(\tau^{\FI}))/d\lambda}
  = \frac{\psi^{\PI}}{\psi^{\FI}} > 1.
\end{equation}

\medskip\noindent\textbf{Step 4: Probability-level amplification.}
In terms of the ruin probabilities themselves, for small changes in~$\lambda$:
\begin{align}\label{eq:prob_sensitivity}
  \frac{d\log\Prob(\tau^{\PI})}{d\lambda}
  &= \psi^{\PI}\cdot a \cdot \frac{1}{1-\lambda}
  \qquad\text{(approximately)}, \\
  \frac{d\log\Prob(\tau^{\FI})}{d\lambda}
  &= 2\gamma\,a. \notag
\end{align}
The ratio of the $\lambda$-elasticities of the ruin probabilities
equals the Lundberg exponent ratio~\eqref{eq:sensitivity_ratio_proof}.
\end{proof}

\begin{remark}[Numerical verification of the sensitivity ratio]
\label{rem:sensitivity_numerical}
Our Monte Carlo simulations (Table~\ref{tab:sensitivity}) yield
empirical Lundberg exponents $\psi_{\mathrm{eff}}^{\PI} \approx 30$--$66$
across different~$\lambda$ values, compared with theoretical
$\psi^{\FI} = 2\gamma(1-\lambda) \in [5.5, 10]$.
The empirical ratio $\psi^{\PI}/\psi^{\FI} \approx 5$--$6$ exceeds
the myopic prediction of~$1.82$, consistent with the additional safety
provided by the hedging term, which further reduces the effective PI
portfolio.
\end{remark}


%% ------------------------------------------------------------------------
\subsection{Proof of Proposition~\ref{prop:heterogeneous}}\label{sec:proof_heterogeneous}
%% ------------------------------------------------------------------------

We outline the perturbation analysis for heterogeneous risk aversion.
Full algebraic details are provided in the supplementary appendix.

\begin{proof}[Proof sketch]
Consider two types of investors: type~A with risk aversion
$\gamma_A = \gamma + \varepsilon$ and type~B with
$\gamma_B = \gamma - \varepsilon$, each comprising $N/2$ of the
$N$~investors.
At $\varepsilon = 0$, this reduces to the symmetric game.

\medskip\noindent\textbf{Step 1: Coupled PDE system.}
Under heterogeneity, the equilibrium functions $g_A$ and $g_B$ satisfy
a coupled Cauchy system
\citep[Theorems~3.1 and 3.2, Remark~3.4]{HuangSun2025}:
\begin{equation}\label{eq:coupled_PDE}
\begin{cases}
  \partial_t c_k + \dfrac{\beta^2}{2}\,\partial_{qq} c_k
  - \dfrac{\beta}{\sigma}\,
    \bigl[\theta - \sigma^2\bar{g}(t,q)\bigr]\,\partial_q c_k
  + \kappa_k\,\dfrac{\theta^2}{\sigma^2} = 0, \\[6pt]
  c_k(T,q) = 0, \qquad k \in \{A, B\},
\end{cases}
\end{equation}
where the aggregate function
$\bar{g} = \frac{1}{N(1-\lambda)}\bigl[\frac{N_A g_A}{\gamma_A}
+ \frac{N_B g_B}{\gamma_B}\bigr]$
couples the two equations through the drift coefficient.

\medskip\noindent\textbf{Step 2: Perturbation expansion.}
Write $c_k = c_0 + \varepsilon\,c_1^k + O(\varepsilon^2)$ and
$g_k = g_0 + \varepsilon\,g_1^k + O(\varepsilon^2)$.
At zeroth order, $c_0$ and $g_0$ satisfy the symmetric Cauchy problem
and are $\lambda$-independent (Remark~\ref{rem:g_indep}).

The first-order corrections satisfy
\begin{align}
  g_1^A &= -\frac{\beta}{\sigma\kappa_0}\,\partial_q c_1^A
          - \frac{\beta\,\partial_q c_0}{\sigma\kappa_0\gamma},
  \label{eq:g1A} \\
  g_1^B &= -\frac{\beta}{\sigma\kappa_0}\,\partial_q c_1^B
          + \frac{\beta\,\partial_q c_0}{\sigma\kappa_0\gamma},
  \label{eq:g1B}
\end{align}
where $\kappa_0 = 1/\gamma$.

\medskip\noindent\textbf{Step 3: Antisymmetry and decoupling.}
With $N_A = N_B = N/2$, the aggregate first-order correction vanishes:
$\bar{g}_1 = (g_1^A + g_1^B)/[2\gamma(1-\lambda)] = 0$, because
$g_1^A = -g_1^B$ (from the antisymmetry $c_1^A = -c_1^B$).
The antisymmetric mode $c_1^D := c_1^A$ satisfies the \emph{decoupled} PDE
\begin{equation}\label{eq:c1D_PDE}
  \mathcal{L}^0 c_1^D + \frac{(1-\lambda/N)\,\theta^2}{\gamma^2\sigma^2} = 0,
  \qquad c_1^D(T,q) = 0,
\end{equation}
where $\mathcal{L}^0$ is the backward Kolmogorov operator of the
zeroth-order filter process.

\medskip\noindent\textbf{Step 4: Stochastic representation and
$\lambda$-dependence.}
By the Feynman--Kac theorem, the solution
of~\eqref{eq:c1D_PDE} has the representation
\begin{equation}\label{eq:c1D_FK}
  c_1^D(t,q) = \frac{(1-\lambda/N)}{\gamma^2\sigma^2}\,
  \E_{Q_0}^{t,q}\!\biggl[\int_t^T \theta\bigl(P_0(u)\bigr)^2\,du\biggr]
  > 0,
\end{equation}
where $P_0$ is the filter process under the zeroth-order
Girsanov measure~$Q_0$.
The source term $\theta^2 > 0$ ensures strict positivity.

The $\lambda$-dependence enters through:
\begin{enumerate}[label=(\alph*)]
\item the explicit prefactor $(1-\lambda/N)$, and
\item the measure~$Q_0$, whose drift coefficient involves
  $\bar{g}_0 = g_0/[\gamma(1-\lambda)]$.
\end{enumerate}
Both channels are absent in the symmetric case, confirming
that heterogeneity breaks the $\lambda$-independence of~$g$.

\medskip\noindent\textbf{Step 5: Effect on the wealth gap.}
The expected wealth gap for type~A at first order in~$\varepsilon$ is
\begin{equation}\label{eq:GA_expansion}
  G_A(\lambda, T) = G_0(\lambda, T)
  - \frac{\varepsilon\mu_1}{\gamma(1-\lambda)}\biggl[
    \frac{I_0^{\mathrm{tot}}}{\gamma} + I_1^{A,\mathrm{tot}}\biggr]
  + O(\varepsilon^2),
\end{equation}
where $I_0^{\mathrm{tot}} = \int_0^T [\mu_1/\sigma^2 -
\E_{\mu_1}[g_0(t,P(t))]]\,dt > 0$ and
$I_1^{A,\mathrm{tot}} = \int_0^T \E_{\mu_1}[g_1^A(t,P(t))]\,dt$.

The $\lambda$-derivative satisfies
\begin{equation}\label{eq:dGA_dlambda}
  \frac{\partial G_A}{\partial\lambda}
  = \frac{G_A}{1-\lambda} + \varepsilon\,\Phi(\lambda, T)
  + O(\varepsilon^2),
\end{equation}
where $\Phi(\lambda, T)$ captures the $\lambda$-correction from
heterogeneity.
The sign of~$\Phi$ is determined by the stochastic representation
of~$c_1^D$: since the source term in~\eqref{eq:c1D_PDE} is proportional
to $(1-\lambda/N)$ and the measure~$Q_0$ makes $\E_{Q_0}[\theta^2]$
increasing in~$\lambda$ (through the drift modification), we obtain
$\Phi(\lambda, T) > 0$ in the mean-field limit $N \to \infty$.
\end{proof}


%% ========================================================================
\section{Numerical Analysis}\label{sec:numerics_v2}
%% ========================================================================

We validate the analytical results of Section~\ref{sec:results} through
extensive Monte Carlo simulations.
All experiments use the parameter calibration of
\citet[Figure~4]{HuangSun2025}, and we simulate both FI and PI strategies
on the \emph{same} Brownian paths to minimize variance in comparisons.

\subsection{Simulation setup}\label{sec:sim_setup}

\paragraph{Base parameters.}
$\mu_1 = 0.2$, $\mu_2 = 0.02$, $\sigma = 0.1$, $\gamma = 5$,
$x = 3$, $p = 0.5$, $T = 10$.
The competition parameter $\lambda$ varies from $0$ to $0.5$.
The warning threshold~$D$ is varied from $1.5$ to $2.95$.

\paragraph{Monte Carlo design.}
We simulate $M = 200{,}000$ paths (increased to $500{,}000$ for
the Lundberg exponent analysis) with time step $\Delta t = 0.01$
using the Euler--Maruyama scheme.
Both FI and PI strategies are simulated using the \emph{same}
Brownian increments $\{\Delta W_k\}_{k=1}^{T/\Delta t}$,
so that the difference $X^{\FI}(t) - X^{\PI}(t)$ has reduced variance.

\paragraph{PI strategy approximation.}
We use the \emph{myopic} PI strategy
$\pi^*(t) = \theta(P(t))/[\sigma^2\gamma(1-\lambda)]$,
omitting the hedging term $h(t,q) \ge 0$.
By the hedging demand analysis in~\eqref{eq:g_decomp}, this overestimates the PI investment
relative to the true equilibrium strategy.
For the good-state comparison
($\Prob_{\mu_1}(\tau^{\PI} \le T) \le \Prob_{\mu_1}(\tau^{\FI} \le T)$),
the myopic approximation is \emph{conservative}: it makes the PI
investor more aggressive than the true equilibrium, bringing PI
closer to FI.
Any safety advantage found for the myopic PI thus also holds for
the true PI.
For the bad-state comparison, the myopic approximation slightly
underestimates the over-investment effect (since the hedging
term further reduces $\pi^*$ below the myopic level).

\paragraph{Simulation under $\Prob_{\mu_k}$.}
To condition on the true state, we set the drift in the stock SDE
to $\mu = \mu_k$ for $k \in \{1, 2\}$.
The filter~$P(t)$ is updated using the Euler--Maruyama discretization
of the innovation process, with the true state determining the
relationship between the stock return and the innovation.


%% ------------------------------------------------------------------------
\subsection{Good-state safety}\label{sec:num_good}
%% ------------------------------------------------------------------------

Table~\ref{tab:good_state} presents the ruin probabilities
$\Prob_{\mu_1}(\tau < T)$ for FI and PI investors under the good state.

\begin{table}[H]
\centering
\caption{Ruin probabilities under $\Prob_{\mu_1}$ (good state).
  Parameters: $\mu_1 = 0.2$, $\mu_2 = 0.02$, $\sigma = 0.1$,
  $\gamma = 5$, $x = 3$, $p = 0.5$, $T = 10$.
  Monte Carlo with $M = 200{,}000$ paths.}
\label{tab:good_state}
\smallskip
\begin{tabular}{ccccl}
\toprule
$D$ & $\lambda$ & $\Prob^{\FI}(\tau < T)$
    & $\Prob^{\PI}(\tau < T)$ & Ratio \\
\midrule
2.0 & 0.0 & 0.000040 & ${<}5{\times}10^{-6}$ & $\approx 0$ \\
2.0 & 0.1 & 0.000115 & ${<}5{\times}10^{-6}$ & $\approx 0$ \\
2.0 & 0.2 & 0.000320 & ${<}5{\times}10^{-6}$ & $\approx 0$ \\
2.0 & 0.3 & 0.000795 & ${<}5{\times}10^{-6}$ & $\approx 0$ \\
\midrule
2.5 & 0.0 & 0.005360 & ${<}5{\times}10^{-6}$ & $\approx 0$ \\
2.5 & 0.1 & 0.008870 & ${<}5{\times}10^{-6}$ & $\approx 0$ \\
2.5 & 0.2 & 0.014700 & ${<}5{\times}10^{-6}$ & $\approx 0$ \\
2.5 & 0.3 & 0.024050 & ${<}5{\times}10^{-6}$ & $\approx 0$ \\
\midrule
2.8 & 0.0 & 0.107305 & ${<}5{\times}10^{-6}$ & $\approx 0$ \\
2.8 & 0.1 & 0.130590 & 0.000010 & $< 0.001$ \\
2.8 & 0.2 & 0.160225 & 0.000055 & $< 0.001$ \\
2.8 & 0.3 & 0.195515 & 0.000230 & $0.001$ \\
\midrule
2.9 & 0.0 & 0.291390 & 0.006685 & $0.023$ \\
2.9 & 0.1 & 0.322010 & 0.014015 & $0.044$ \\
2.9 & 0.2 & 0.355450 & 0.027360 & $0.077$ \\
2.9 & 0.3 & 0.393355 & 0.049635 & $0.126$ \\
\bottomrule
\end{tabular}
\end{table}

\emph{Discussion.}
The PI investor is dramatically safer than the FI investor under the
good state, with zero observed ruin for $D \le 2.8$ across all~$\lambda$
values in $200{,}000$ simulated paths.
Even at $D = 2.9$ (safety margin $a = 0.1$), the PI ruin probability is
2\%--13\% of the FI level depending on~$\lambda$.
This is consistent with Theorem~\ref{thm:expo_bound}:
the PI ruin probability is bounded by the FI \emph{infinite-horizon}
ruin probability, which for $a \ge 0.2$ is vanishingly small.

The mechanism is transparent: the PI investor starts at
$\pi^*(0) = 2.20$ (at $\lambda = 0$) compared to $\pi^{\FI} = 4.00$,
giving 45\% less leverage.
This reduced leverage during the early learning phase---when the filter
is still uncertain---provides an exponential safety buffer.

The ratio $\Prob^{\PI}/\Prob^{\FI}$ increases with~$\lambda$
(from $0.023$ to $0.126$ at $D = 2.9$), indicating that competition
erodes PI's safety advantage, but never eliminates it within the
parameter range studied.


%% ------------------------------------------------------------------------
\subsection{Bad-state amplification}\label{sec:num_bad}
%% ------------------------------------------------------------------------

Table~\ref{tab:bad_state} presents the ruin probabilities under the bad
state $\Prob_{\mu_2}$.

\begin{table}[H]
\centering
\caption{Ruin probabilities under $\Prob_{\mu_2}$ (bad state).
  Same parameters as Table~\ref{tab:good_state}.}
\label{tab:bad_state}
\smallskip
\begin{tabular}{ccccl}
\toprule
$D$ & $\lambda$ & $\Prob^{\FI}(\tau < T)$
    & $\Prob^{\PI}(\tau < T)$ & Ratio \\
\midrule
2.0 & 0.0 & ${<}5{\times}10^{-6}$ & ${<}5{\times}10^{-6}$ & --- \\
2.0 & 0.1 & ${<}5{\times}10^{-6}$ & ${<}5{\times}10^{-6}$ & --- \\
2.0 & 0.2 & ${<}5{\times}10^{-6}$ & ${<}5{\times}10^{-6}$ & --- \\
2.0 & 0.3 & ${<}5{\times}10^{-6}$ & ${<}5{\times}10^{-6}$ & --- \\
\midrule
2.5 & 0.0 & 0.000010 & 0.000010 & $1.00$ \\
2.5 & 0.1 & 0.000035 & 0.000090 & $2.57$ \\
2.5 & 0.2 & 0.000135 & 0.000415 & $3.07$ \\
2.5 & 0.3 & 0.000770 & 0.001640 & $2.13$ \\
\midrule
2.8 & 0.0 & 0.034125 & 0.039925 & $1.17$ \\
2.8 & 0.1 & 0.052340 & 0.056370 & $1.08$ \\
2.8 & 0.2 & 0.078290 & 0.078575 & $1.00$ \\
2.8 & 0.3 & 0.114105 & 0.109160 & $0.96$ \\
\midrule
2.9 & 0.0 & 0.228810 & 0.218125 & $0.95$ \\
2.9 & 0.1 & 0.268700 & 0.259975 & $0.97$ \\
2.9 & 0.2 & 0.314375 & 0.310615 & $0.99$ \\
2.9 & 0.3 & 0.367460 & 0.365875 & $1.00$ \\
\bottomrule
\end{tabular}
\end{table}

\emph{Discussion.}
The risk reversal phenomenon is clearly visible.
At moderate safety margin ($D = 2.5$, $a = 0.5$) with
$\lambda \in \{0.1, 0.2, 0.3\}$, the PI ruin probability exceeds the
FI one by a factor of $2$--$3$, confirming the bad-state amplification
of Theorem~\ref{thm:bad_state}.
The reversal is driven by the over-investment mechanism of
Proposition~\ref{prop:over_invest}: the PI investor's initial portfolio
$\bar{\pi} = 3.14$ (at $\lambda = 0.3$) is $5.5\times$ the FI optimal
$\pi_2^{\FI} = 0.57$.

As $D$ increases toward~$x$ (small safety margin), the bad-state
ratio approaches~$1$ from below for $D = 2.8$ at large~$\lambda$ and
$D = 2.9$ across all~$\lambda$.
This is because with a tiny safety margin, the transient
over-investment effect is overwhelmed by the symmetric
barrier-hitting behavior at short times.

\medskip
Table~\ref{tab:mixture} presents the ruin probabilities under the
mixture measure $\Prob_{p} = p\,\Prob_{\mu_1} + (1-p)\,\Prob_{\mu_2}$,
which represents the investor's \emph{ex ante} risk.

\begin{table}[H]
\centering
\caption{Ruin probabilities under the mixture $\Prob_{0.5}$.}
\label{tab:mixture}
\smallskip
\begin{tabular}{ccccl}
\toprule
$D$ & $\lambda$ & $\Prob_{\mathrm{mix}}^{\FI}$ & $\Prob_{\mathrm{mix}}^{\PI}$ & Ratio \\
\midrule
2.0 & 0.0 & 0.000020 & ${<}5{\times}10^{-6}$ & $\approx 0$ \\
2.0 & 0.3 & 0.000398 & ${<}5{\times}10^{-6}$ & $\approx 0$ \\
\midrule
2.5 & 0.0 & 0.002685 & 0.000005 & $0.002$ \\
2.5 & 0.3 & 0.012410 & 0.000820 & $0.066$ \\
\midrule
2.8 & 0.0 & 0.070715 & 0.019963 & $0.282$ \\
2.8 & 0.3 & 0.154810 & 0.054695 & $0.353$ \\
\midrule
2.9 & 0.0 & 0.260100 & 0.112405 & $0.432$ \\
2.9 & 0.3 & 0.380408 & 0.207755 & $0.546$ \\
\bottomrule
\end{tabular}
\end{table}

\emph{Discussion.}
Under the mixture measure (which reflects the investor's prior
uncertainty), the PI investor is safer overall: the ratio
$\Prob_{\mathrm{mix}}^{\PI}/\Prob_{\mathrm{mix}}^{\FI} < 1$
everywhere.
The good-state safety advantage dominates the bad-state disadvantage
because the bad-state ruin events are rare in absolute terms.
However, the ratio increases with both~$D$ and~$\lambda$,
approaching~$1$ when the safety margin is small and competition
is intense.


%% ------------------------------------------------------------------------
\subsection{Sensitivity analysis}\label{sec:num_sensitivity}
%% ------------------------------------------------------------------------

We investigate the $\lambda$-sensitivity of the ruin probability under
the good state, as predicted by Theorem~\ref{thm:sensitivity}.

\begin{table}[H]
\centering
\caption{Log-ruin probability $-\log\Prob(\tau < T)$ versus~$\lambda$
  under $\Prob_{\mu_1}$, at $D = 2.9$.}
\label{tab:sensitivity}
\smallskip
\begin{tabular}{cccl}
\toprule
$\lambda$ & $-\log\Prob^{\FI}$ & $-\log\Prob^{\PI}$ & Log-ratio \\
\midrule
0.00 & 1.000 & 4.956 & $-3.956$ \\
0.05 & 0.950 & 4.740 & $-3.790$ \\
0.10 & 0.900 & 4.233 & $-3.333$ \\
0.15 & 0.850 & 3.917 & $-3.067$ \\
0.20 & 0.800 & 3.569 & $-2.769$ \\
0.25 & 0.750 & 3.268 & $-2.518$ \\
0.30 & 0.700 & 2.977 & $-2.277$ \\
0.35 & 0.650 & 2.721 & $-2.071$ \\
0.40 & 0.600 & 2.465 & $-1.865$ \\
0.45 & 0.550 & 2.246 & $-1.696$ \\
0.50 & 0.500 & 2.009 & $-1.509$ \\
\bottomrule
\end{tabular}
\end{table}

\emph{Discussion.}
Both $-\log\Prob^{\FI}$ and $-\log\Prob^{\PI}$ decrease linearly
in~$\lambda$ (ruin probability increases with competition), but the
PI slope is approximately $5$--$6$ times steeper.
Specifically:
\begin{itemize}
\item FI: $d(-\log\Prob^{\FI})/d\lambda \approx -1.0$
  (from $\psi^{\FI} = 2\gamma(1-\lambda) \cdot a$ with $a = 0.1$).
\item PI: $d(-\log\Prob^{\PI})/d\lambda \approx -5.9$
  (average slope from Table~\ref{tab:sensitivity}).
\item Ratio: $\approx 5.9$, consistent with the empirical Lundberg
  exponent ratio $\psi^{\PI}/\psi^{\FI} \approx 5$--$6$.
\end{itemize}
This confirms the sensitivity amplification of
Theorem~\ref{thm:sensitivity}: the same increase in competition
intensity~$\lambda$ has a disproportionately larger effect on the
PI ruin probability.

Figure~\ref{fig:logP_vs_lambda} visualizes this relationship.

\begin{figure}[H]
\centering
\includegraphics[width=0.7\textwidth]{figures/sensitivity_logP_vs_lambda.png}
\caption{Log-ruin probability $\log\Prob(\tau < T)$ versus competition
  intensity~$\lambda$ at $D = 2.9$.
  The PI curve (solid) has a steeper slope than FI (dashed),
  demonstrating the sensitivity amplification of
  Theorem~\ref{thm:sensitivity}.}
\label{fig:logP_vs_lambda}
\end{figure}


\subsection{Effective Lundberg exponents}\label{sec:num_lundberg}

We estimate the effective Lundberg exponent from the dependence of
the ruin probability on the safety margin~$a = x - D$.
The theoretical prediction is $\psi^{\FI} = 2\gamma(1-\lambda)$
for FI; for PI, we estimate $\psi^{\PI}$ from the log-linear
fit $-\log\Prob(\tau < T) \approx \psi \cdot a$ across
multiple~$D$ values, using $T = 10$ as a proxy for $T = \infty$.

\begin{table}[H]
\centering
\caption{Effective Lundberg exponents at $T = 10$.
  Monte Carlo with $M = 500{,}000$ paths.}
\label{tab:lundberg}
\smallskip
\begin{tabular}{cccc}
\toprule
$\lambda$ & $\psi^{\FI}$ (theory)
          & $\psi^{\PI}_{\mathrm{eff}}$ (empirical) & Ratio \\
\midrule
0.00 & 10.00 & 57.76 & 5.78 \\
0.15 &  8.50 & 48.43 & 5.70 \\
0.30 &  7.00 & 35.78 & 5.11 \\
0.45 &  5.50 & 28.32 & 5.15 \\
\bottomrule
\end{tabular}
\end{table}

\emph{Discussion.}
The empirical Lundberg exponents confirm that the PI ruin probability
decays exponentially faster with the safety margin than the FI one,
with a ratio of approximately $5$--$6$.
This ratio exceeds the myopic prediction
$\mu_1/\theta(p) = 0.2/0.11 \approx 1.82$ by a factor of~$3$,
demonstrating that the hedging term provides additional safety
beyond the myopic under-investment.

The ratio is remarkably stable across~$\lambda$ values
($5.11$--$5.78$), consistent with the theoretical prediction that
$\psi^{\PI}/\psi^{\FI}$ depends on the filter dynamics (which are
$\lambda$-independent) rather than on the competition structure.

Figure~\ref{fig:lundberg} illustrates the exponential decay of
ruin probability with the safety margin.

\begin{figure}[H]
\centering
\includegraphics[width=0.7\textwidth]{figures/lundberg_exponent.png}
\caption{Effective Lundberg exponents.
  Log-ruin probability $\log\Prob(\tau < T)$ versus safety margin~$a$
  for FI (dashed) and PI (solid) at several~$\lambda$ values.
  The steeper PI slopes reflect the larger Lundberg exponent.}
\label{fig:lundberg}
\end{figure}


%% ------------------------------------------------------------------------
\subsection{Boundary behavior}\label{sec:num_boundary}
%% ------------------------------------------------------------------------

We examine the behavior of the PI/FI ruin ratio as the barrier~$D$
approaches the initial wealth~$x$, and as $D \to 0$.

\begin{table}[H]
\centering
\caption{PI/FI ruin ratio near the boundary ($D \to x$) under $\Prob_{\mu_1}$.}
\label{tab:boundary}
\smallskip
\begin{tabular}{cccccl}
\toprule
$D$ & $a$ & $\lambda$ & $\Prob^{\FI}$ & $\Prob^{\PI}$ & Ratio \\
\midrule
2.0 & 1.00 & 0.0 & 0.000045 & ${<}5{\times}10^{-6}$ & $\approx 0$ \\
2.0 & 1.00 & 0.4 & 0.002479 & ${<}5{\times}10^{-6}$ & $\approx 0$ \\
\midrule
2.6 & 0.40 & 0.0 & 0.018316 & ${<}5{\times}10^{-6}$ & $\approx 0$ \\
2.6 & 0.40 & 0.4 & 0.090718 & ${<}5{\times}10^{-6}$ & $\approx 0$ \\
\midrule
2.8 & 0.20 & 0.0 & 0.135335 & ${<}5{\times}10^{-6}$ & $\approx 0$ \\
2.8 & 0.20 & 0.4 & 0.301194 & 0.001380 & $0.005$ \\
\midrule
2.9 & 0.10 & 0.0 & 0.367879 & 0.007180 & $0.020$ \\
2.9 & 0.10 & 0.4 & 0.548812 & 0.083805 & $0.153$ \\
\midrule
2.95 & 0.05 & 0.0 & 0.606531 & 0.134955 & $0.222$ \\
2.95 & 0.05 & 0.4 & 0.740818 & 0.295725 & $0.399$ \\
\bottomrule
\end{tabular}
\end{table}

\emph{Discussion.}
As $D \to x$ (i.e., $a \to 0$), both ruin probabilities approach~$1$
(any process will cross a nearby barrier), and the PI/FI ratio increases
toward~$1$.
However, the ratio never exceeds~$1$ within the observed range:
PI remains safer everywhere in $(D, \lambda)$-space under
$\Prob_{\mu_1}$.

Figure~\ref{fig:boundary} shows the ratio as a function of~$D$.

\begin{figure}[H]
\centering
\includegraphics[width=0.7\textwidth]{figures/boundary_ratio_vs_D.png}
\caption{PI/FI ruin ratio versus barrier level~$D$ under $\Prob_{\mu_1}$
  for several~$\lambda$ values.
  The ratio is identically zero for $D \le 2.6$ (PI has zero observed ruin)
  and increases toward~$1$ as $D \to x = 3$, but never crosses~$1$.}
\label{fig:boundary}
\end{figure}

As $D \to 0$, both ruin probabilities vanish exponentially fast (with
rates $\psi^{\FI}\cdot a$ and $\psi^{\PI}\cdot a$), and the ratio
$\Prob^{\PI}/\Prob^{\FI} \to 0$ since $\psi^{\PI} > \psi^{\FI}$.

As $T \to \infty$, the Lundberg exponents dominate and the ratio
converges to $\exp(-(\psi^{\PI} - \psi^{\FI})\cdot a)$, which
provides an exponentially small upper bound on the long-run PI/FI ratio.


\subsection{Phase diagram}\label{sec:num_phase}

Figure~\ref{fig:phase} presents the heatmap of the PI/FI ruin ratio
in $(D, \lambda)$-space under~$\Prob_{\mu_1}$.

\begin{figure}[H]
\centering
\includegraphics[width=0.7\textwidth]{figures/phase_diagram.png}
\caption{Phase diagram of the PI/FI ruin ratio in $(D, \lambda)$-space
  under $\Prob_{\mu_1}$.
  Blue (ratio $\approx 0$): PI has negligible ruin.
  The ratio increases toward the upper-right corner
  ($D \to x$, $\lambda \to 1$) but remains below~$1$ everywhere.}
\label{fig:phase}
\end{figure}

\emph{Discussion.}
The phase diagram confirms the uniform safety of PI under the good state
across the entire $(D, \lambda)$-space.
The ratio is smallest (PI most advantaged) at moderate~$D$ and
small~$\lambda$, where the FI ruin probability is already non-negligible
but the PI under-investment provides strong protection.
The PI advantage is greatest precisely in the regime where ruin risk
is most economically relevant.


\subsection{Robustness: alternative parameter calibrations}
\label{sec:num_robust}

To verify that our findings are not artifacts of the specific
Huang--Sun calibration, we repeat the two-state analysis for
two additional parameter sets.

\begin{table}[H]
\centering
\caption{Risk reversal under alternative calibrations.
  Panel~B: $\gamma = 2$, $\sigma = 0.2$ (moderate).
  Panel~C: $\gamma = 1$, $\sigma = 0.1$ (aggressive).}
\label{tab:robust}
\smallskip
\begin{tabular}{lcccc}
\toprule
Panel & $D$ & $\lambda$ & Good: PI/FI & Bad: PI/FI \\
\midrule
B ($\gamma{=}2$, $\sigma{=}0.2$)
  & 2.0 & 0.3 & 0.000 & 24.0 \\
B & 2.5 & 0.1 & 0.000 & 10.6 \\
B & 2.5 & 0.3 & 0.002 & 5.1 \\
\midrule
C ($\gamma{=}1$, $\sigma{=}0.1$)
  & 1.0 & 0.0 & 0.000 & 3.1 \\
C & 1.5 & 0.3 & 0.000 & 1.2 \\
C & 2.0 & 0.0 & 0.000 & 1.2 \\
\bottomrule
\end{tabular}
\end{table}

\emph{Discussion.}
The risk reversal phenomenon persists across all parameter
calibrations: PI is safer under the good state (ratio $< 1$) and
riskier under the bad state (ratio $> 1$) for moderate safety margins
and competition intensity.
Under more aggressive calibrations (Panel~C, $\gamma = 1$), the
reversal amplitudes are smaller because the FI bad-state portfolio
$\pi_2^{\FI} = 2.0$ is already substantial, narrowing the gap with
the PI portfolio.


%% ========================================================================
\section{Discussion}\label{sec:discuss_v2}
%% ========================================================================

\subsection{Economic interpretation of the risk reversal}
\label{sec:discuss_reversal}

The central finding of this paper---that partial information
\emph{protects} investors under the good state but \emph{endangers}
them under the bad state---admits a transparent economic interpretation.

\paragraph{The paradox of ignorance.}
Under the good state ($\mu = \mu_1$), the FI investor knows the true
(high) drift and invests aggressively:
$\pi^{\FI} = \mu_1/[\sigma^2\gamma(1-\lambda)]$.
The PI investor, uncertain about the drift, uses the estimated
value $\theta(P(t)) < \mu_1$ and invests more conservatively.
This under-investment is suboptimal in expectation
(Theorem~\ref{thm:wealth} shows a wealth loss), but it provides
an unintended benefit: reduced portfolio volatility.

The reduced volatility translates directly into reduced ruin
probability.
The PI investor's wealth process has instantaneous variance
$(\pi^*(t))^2\sigma^2 < (\pi^{\FI})^2\sigma^2$ during the early
learning phase, which is precisely when the safety margin is
smallest and ruin risk is highest.
By the time the filter converges ($P(t) \approx 1$) and the PI
portfolio approaches the FI level, the accumulated drift has
already moved wealth far from the barrier.

This is a \emph{mechanical} effect, not a reflection of superior
judgment: the PI investor is simply too cautious to be endangered.
The protection is strongest when uncertainty is highest (large
$\delta$, small~$p$) and the barrier is moderately close
($D \in [2.0, 2.8]$).

\paragraph{Under the bad state: the price of ignorance.}
The mirror image applies under $\mu = \mu_2$.
The FI investor knows the low drift and invests conservatively:
$\pi_2^{\FI} = \mu_2/[\sigma^2\gamma(1-\lambda)]$, which is
$5.5\times$ smaller than the PI investor's initial position.
The PI investor, still believing $\mu$ could be $\mu_1$ with
probability~$p = 0.5$, over-invests dramatically.
This over-investment generates unnecessary volatility and
increases ruin risk, especially during the early period before
the filter can identify the bad state.

\paragraph{The aggregate picture.}
Under the ex ante mixture measure $\Prob_p$
($p = 0.5$), the PI investor is safer overall
(Table~\ref{tab:mixture}).
The good-state safety advantage dominates because:
(i)~the PI under-investment is symmetric (about $50\%$ of FI)
in the good state, while the over-investment under the bad state,
though large in ratio, is small in absolute portfolio size;
(ii)~the exponential dependence on the Lundberg exponent means
that the good-state advantage ($\psi^{\PI}/\psi^{\FI} \approx 5$)
amplifies exponentially with the safety margin.


%% ------------------------------------------------------------------------
\subsection{Connection to Huang--Sun's amplification observation}
\label{sec:discuss_HS}
%% ------------------------------------------------------------------------

\citet[Figure~4]{HuangSun2025} observe that co-drawdown probabilities
increase with~$\lambda$ more sharply under PI than under FI, and
conjecture that partial information amplifies downward risk.
Our analysis provides a precise theoretical explanation.

The amplification is not about the \emph{level} of ruin probability
(PI can be safer or riskier depending on the true state), but about
the \emph{sensitivity} to competition.
Theorem~\ref{thm:sensitivity} shows that the Lundberg exponent is
$\psi^{\PI}/\psi^{\FI} \approx 5$--$6$ times larger for PI, meaning
that the same increase in~$\lambda$ causes a proportionally larger
change in $\log\Prob(\tau^{\PI})$.

Under the mixture measure (Figure~4 of \citealt{HuangSun2025}),
the increasing co-drawdown reflects the bad-state channel:
as~$\lambda$ increases, the bad-state PI ruin probability grows
rapidly (because the over-investment is amplified by the
$1/(1-\lambda)$ factor), while the good-state PI ruin remains
near zero.
The mixture probability thus shifts weight toward the bad-state
contribution, generating the observed upward trend.


%% ------------------------------------------------------------------------
\subsection{Implications for systemic risk}\label{sec:discuss_systemic}
%% ------------------------------------------------------------------------

The risk reversal has important implications for the macroprudential
assessment of fund manager competition under incomplete information.

\paragraph{Tail amplification.}
Under the mixture measure, the PI investor is safer \emph{on average}
(Table~\ref{tab:mixture}).
However, Theorem~\ref{thm:bad_state} shows that conditional on the
bad state, PI investors face catastrophically higher ruin risk.
This is a form of \emph{tail amplification}: partial information
shifts risk from the body to the tail of the state distribution.

A regulator using the mixture measure would underestimate the
systemic risk posed by PI investors, because the good-state safety
masks the bad-state vulnerability.
The correct assessment requires conditioning on the worst case,
which is precisely the scenario where PI investors are most
vulnerable.

\paragraph{Simultaneous failure.}
Under the symmetric game, all $N$~PI investors use the same strategy
(driven by the common filter) and therefore fail simultaneously:
$\tau_i = \tau_j$ a.s.\ for all~$i, j$.
This perfect correlation in default events means that the co-drawdown
probability equals the individual ruin probability.
The heterogeneous game (Proposition~\ref{prop:heterogeneous}) introduces
some dispersion in failure times, but the common filter ensures
that the correlation remains high.

\paragraph{Policy implications.}
Our results suggest that:
\begin{enumerate}[label=(\roman*)]
\item Competition ($\lambda$) is more dangerous under partial
  information than full information, not because PI investors are
  riskier on average, but because they are more \emph{sensitive}
  to competition.
\item Disclosure requirements that improve information quality
  (moving the regime from PI toward FI) have a dual effect:
  they increase average expected wealth
  (Theorem~\ref{thm:wealth})
  but may also increase ruin risk under the good state
  (Theorem~\ref{thm:expo_bound}).
\item The optimal regulatory response depends on the prior
  probability~$p$ of the good state: when $p$ is high
  (economy likely in good shape), the good-state safety of PI
  may be more relevant than the bad-state vulnerability.
\end{enumerate}


%% ------------------------------------------------------------------------
\subsection{Open problems}\label{sec:open}
%% ------------------------------------------------------------------------

\begin{enumerate}[label=\arabic*.]
\item \textbf{Full pathwise comparison under $\Prob_{\mu_1}$.}
Theorem~\ref{thm:expo_bound} establishes
$\Prob_{\mu_1}(\tau^{\PI} \le T) \le \Prob_{\mu_1}(\tau^{\FI} < \infty)$
but falls short of the full comparison
$\Prob_{\mu_1}(\tau^{\PI} \le T) \le \Prob_{\mu_1}(\tau^{\FI} \le T)$.
The obstruction is a Girsanov density correlation: under the true
measure~$\Prob_{\mu_1}$, the Dubins--Schwarz comparison established
under~$Q$ does not transfer because the Girsanov density~$L(T)$
correlates differently with the two ruin events.
Resolving this likely requires either a direct coupling argument
on the wealth processes (not on the time-changed Brownian motions)
or a comparison theorem for one-dimensional diffusions with
stochastic coefficients.

\item \textbf{Sharp boundary of the bad-state reversal region.}
Our numerics (Table~\ref{tab:bad_state}) show that the
bad-state amplification $\Prob^{\PI}/\Prob^{\FI} > 1$ occurs
for moderate~$D$ and~$\lambda$, but the precise boundary in
$(D, \lambda)$-space is determined only numerically.
A sharp characterization---perhaps in terms of the over-investment
ratio $\bar{\pi}/\pi_2^{\FI}$ and the BL formula---would provide
closed-form conditions for the reversal.

\item \textbf{General $N$-type heterogeneous game.}
Proposition~\ref{prop:heterogeneous} treats two types with equal population.
The general case with $K$~types and arbitrary population fractions
$\alpha_1, \ldots, \alpha_K$ involves a $K$-dimensional coupled PDE
system whose perturbation theory is substantially more complex.
Of particular interest is whether the antisymmetry
($c_1^A = -c_1^B$) persists for unequal populations, and whether
the aggregate correction $\bar{g}_1$ remains zero.

\item \textbf{Price impact and endogenous volatility.}
The current model takes the stock price as exogenous.
In a market with price impact, the aggregate portfolio
$\bar{\pi}$ feeds back into the stock dynamics, creating a
fixed-point problem.
Under PI, the filter itself depends on the stock price, so the
feedback loop involves both the portfolio and the belief dynamics.
The interaction between learning, price impact, and competition
is a challenging but important direction.

\item \textbf{Mean-field game limit.}
As $N \to \infty$ with $\lambda$ fixed, the game converges to a
mean-field game.
The convergence rate and the behavior of the ruin probability in
the large-$N$ limit are open questions.
In particular, the heterogeneous perturbation
(Proposition~\ref{prop:heterogeneous}) suggests that the mean-field limit
may qualitatively differ from the finite-$N$ game when
risk-aversion heterogeneity is present.
\end{enumerate}


%% ========================================================================
\section{Conclusion}\label{sec:conclude_v2}
%% ========================================================================

We have established a theoretical framework for first-passage time
comparison in mean-variance portfolio games under partial information,
building on the model of \citet{HuangSun2025}.
Our main contribution is the \emph{risk reversal theorem}: the PI
investor is dramatically safer under the good state
(Theorem~\ref{thm:expo_bound}) but faces amplified ruin risk under
the bad state (Theorem~\ref{thm:bad_state}).

The good-state safety arises from the mechanical under-investment of
the uncertain PI investor.
The exponential martingale bound
(Theorem~\ref{thm:expo_bound}) shows that the PI ruin probability
is bounded by the FI \emph{infinite-horizon} probability, with a
theoretical Lundberg exponent ratio of $\mu_1/\theta(p) \approx 1.82$
(Monte Carlo estimates: $5$--$6$ times).
The bad-state amplification arises from the over-investment of the
PI investor who does not know the state is bad.
The coupling argument
(Theorem~\ref{thm:bad_state}) provides an explicit lower bound on
the PI ruin probability in terms of the Bachelier--L\'evy formula
and the over-investment ratio $\bar{\pi}/\pi_2^{\FI}$.

The sensitivity amplification (Theorem~\ref{thm:sensitivity}) provides
the theoretical foundation for the co-drawdown phenomenon observed in
\citet[Figure~4]{HuangSun2025}: the PI ruin probability is
$5$--$6$ times more sensitive to the competition parameter~$\lambda$,
explaining why co-drawdown increases sharply with competition under
partial information.

The perturbation analysis for heterogeneous investors
(Proposition~\ref{prop:heterogeneous}) reveals that risk-aversion
heterogeneity breaks the $\lambda$-independence of the equilibrium
function~$g$, creating a new channel through which competition affects
the filter dynamics.
This opens the door to a richer theory of co-drawdown in heterogeneous
populations, where different investor types fail at different rates
and the aggregate default behavior depends on the composition of the
investor population.


%% ========================================================================
%% Bibliography
%% ========================================================================

\bibliographystyle{plainnat}

\begin{thebibliography}{99}


\bibitem[Basak and Makarov(2014)]{BasakMakarov2014}
Basak, S. and Makarov, D. (2014).
\newblock Strategic asset allocation in money management.
\newblock \emph{Journal of Finance}, 69(1):179--217.

\bibitem[B{\"a}uerle and Bayraktar(2013)]{BaeuerleBayraktar2013}
B{\"a}uerle, N. and Bayraktar, E. (2013).
\newblock A note on applications of stochastic ordering to control problems
  in insurance and finance.
\newblock \emph{Stochastics}, 86(2):330--340.
\newblock arXiv:1312.1068.

\bibitem[Carmona et~al.(2018)]{Carmona2018}
Carmona, R., Fouque, J.-P., Mousavi, S.~M., and Sun, L.-H. (2018).
\newblock Systemic risk and stochastic games with delay.
\newblock \emph{Journal of Optimization Theory and Applications},
  179(1):88--120.

\bibitem[Deng et~al.(2022)]{Deng2022}
Deng, C., Su, X., and Yang, C. (2022).
\newblock Relative wealth concerns with partial information and heterogeneous
  priors.
\newblock \emph{SIAM Journal on Financial Mathematics}, 13(3):SC59--SC85.
\newblock arXiv:2007.11781.

\bibitem[Espinosa and Touzi(2015)]{EspinosaTouzi2015}
Espinosa, G.-E. and Touzi, N. (2015).
\newblock Optimal investment under relative performance concerns.
\newblock \emph{Mathematical Finance}, 25(2):221--257.

\bibitem[Henderson et~al.(2020)]{Henderson2020}
Henderson, V., Kladivko, K., Monoyios, M., and Reisinger, C. (2020).
\newblock Executive stock option exercise with full and partial information
  on a drift change point.
\newblock \emph{SIAM Journal on Financial Mathematics}, 11(4):1007--1062.

\bibitem[Huang and Sun(2025)]{HuangSun2025}
Huang, Y.-J. and Sun, Z. (2025).
\newblock Mean-variance portfolio games under partial information.
\newblock \emph{Mathematical Finance}, to appear.


\bibitem[Karatzas and Shreve(1991)]{KaratzasShreve1991}
Karatzas, I. and Shreve, S.~E. (1991).
\newblock \emph{Brownian Motion and Stochastic Calculus}.
\newblock Springer, 2nd edition.

\bibitem[Lacker and Zariphopoulou(2019)]{LackerZariphopoulou2019}
Lacker, D. and Zariphopoulou, T. (2019).
\newblock Mean field and $n$-agent games for optimal investment under
  relative performance criteria.
\newblock \emph{Annals of Applied Probability}, 29(2):1003--1038.

\bibitem[Leobacher and Ngare(2016)]{Leobacher2016}
Leobacher, G. and Ngare, P. (2016).
\newblock Bayesian dividend optimization and finite time ruin probabilities.
\newblock \emph{Stochastic Models}, 32(2):216--242.


\bibitem[Zhang and Yan(2024)]{ZhangYan2024}
Zhang, P. and Yan, Z. (2024).
\newblock Optimal portfolio with relative performance and CRRA risk preferences
  in a partially observable financial market.
\newblock \emph{Applied Mathematics and Computation}, 481:128935.

\end{thebibliography}

\end{document}
